\documentclass[10pt]{article}
\usepackage[a4paper,margin=20mm]{geometry}
\usepackage{fontspec}
\usepackage{xeCJK}
\usepackage[table]{xcolor}
\usepackage{booktabs}
\usepackage{longtable}
\usepackage{tabularx}
\usepackage{array}
\usepackage{enumitem}
\usepackage{amsmath}
\usepackage{hyperref}
\usepackage{fancyhdr}

\setmainfont{TeX Gyre Termes}
\setsansfont{TeX Gyre Heros}
\setCJKmainfont{Droid Sans Fallback}
\setCJKsansfont{Droid Sans Fallback}
\setCJKmonofont{Droid Sans Fallback}
\hypersetup{colorlinks=false,pdfborder={0 0 0}}

\definecolor{titleblue}{RGB}{22,91,125}
\definecolor{softrow}{RGB}{232,240,244}
\definecolor{softblue}{RGB}{235,244,249}

\pagestyle{fancy}
\fancyhf{}
\fancyhead[L]{\sffamily\footnotesize SI / RSI CUDA 并行与优化说明}
\fancyhead[R]{\sffamily\footnotesize 当前实现}
\fancyfoot[C]{\sffamily\footnotesize \thepage}
\renewcommand{\headrulewidth}{0.35pt}

\setlength{\parindent}{0pt}
\setlength{\parskip}{5pt}
\setlist[itemize]{leftmargin=1.5em,itemsep=2pt,topsep=2pt}
\setlist[enumerate]{leftmargin=1.7em,itemsep=2pt,topsep=2pt}

\newcommand{\sect}[1]{\vspace{3mm}{\sffamily\Large\bfseries\color{titleblue} #1}\par\vspace{1mm}{\color{titleblue}\rule{\linewidth}{0.45pt}}\vspace{1mm}}
\newcommand{\subsect}[1]{\vspace{2mm}{\sffamily\large\bfseries #1}\par}
\newcommand{\code}[1]{\texttt{#1}}

\begin{document}

{\sffamily\bfseries\fontsize{24}{28}\selectfont\color{titleblue} SI / RSI CUDA 并行与优化逻辑说明}\par
\vspace{2mm}
{\sffamily\large 当前 Figure 5 实现：S4 coarse SI，S128 fine SI / RSI / RSI-tail，8192 样本}\par
\vspace{5mm}

\begin{tabularx}{\linewidth}{>{\bfseries}l X}
\toprule
对象 & 当前说明\\
\midrule
空间网格 & 30k: 36,470 个单元；230k: 230,111 个单元。\\
角度离散 & coarse SI 使用 $S_4$，$M=24$；fine SI、RSI、RSI-tail 使用 $S_{128}$，$M=16640$。\\
随机样本 & RSI 使用 8192 条独立样本链；RSI-tail 使用同一批链的 $N$ 到 $N+10$ 步平均。\\
CUDA 范围 & 当前 GPU 路径支持各向同性散射、$G=1$；Figure 5 的 SI、RSI、RSI-tail 均走 CUDA。\\
主要代码入口 & Figure 5 的 GPU 入口在 \code{RSISolver::runFigure5GPU()}；大角度/大网格自动进入流式 plan 路径。\\
\bottomrule
\end{tabularx}

\sect{1. 总体计算流程}

Figure 5 中有两类确定性/随机输运计算：

\begin{itemize}
\item SI 是确定性角度求积。每轮迭代中，对所有角方向做扫掠，得到角通量，再按权重累加成标量通量 $\phi$。迭代到相邻两轮的相对 $L_2$ 范数小于阈值。
\item RSI 是随机方向链。每个样本在每一步随机选一个角方向并扫掠，最终对所有样本的第 $N$ 步结果求均值。RSI-tail 不重新跑链，而是在同一条链上累计第 $N$ 到 $N+10$ 步。
\end{itemize}

当前 GPU 路径先上传静态问题数据，包括 mesh、角方向、权重和边界源参数。然后：

\begin{enumerate}
\item 先计算 coarse SI。coarse 角度数小，直接使用 SI-only CUDA API。
\item 再计算 fine SI。对 $S_{128}$ 的 16640 个方向按方向块流式处理 sweep plan。
\item fine SI 收敛得到 $N$ 后，生成 RSI 的随机方向调度表。
\item RSI 按样本批执行。每个样本批根据实际用到的去重方向构造紧凑 plan，在 GPU 上推进所有随机链。
\item 最后只把 SI、RSI、RSI-tail 三个最终单元中心场拷回 CPU 写 CSV。
\end{enumerate}

\sect{2. Sweep Plan 的作用}

非结构四面体网格的扫掠不是简单的规则网格扫描。每个角方向 $\Omega_m$ 都需要一个单元顺序：

\begin{itemize}
\item 对无环方向，单元依赖图可以拓扑排序。代码保存 \code{order} 和 \code{levelOffsets}。
\item 同一个依赖层内的单元没有上风依赖，可以并行计算。
\item 对有环方向，使用投影排序加有限次 Gauss-Seidel pass 作为回退路径。
\end{itemize}

因此 sweep plan 是 GPU 并行的关键中间结构。它包括：

\begin{itemize}
\item \code{orders}: 每个方向的单元访问顺序。
\item \code{levelOffsets}: 无环方向的波前依赖层边界。
\item[] \vspace{-1mm}
\colorbox{softblue}{\parbox{0.95\linewidth}{\footnotesize
例如 \code{levelCells = [5, 9, 12, 3, 7, 8, 20]}，\code{levelOffsets = [0, 3, 6, 7]}，
表示第 0 层是 \code{levelCells[0:3] = [5, 9, 12]}，第 1 层是 \code{levelCells[3:6] = [3, 7, 8]}，
第 2 层是 \code{levelCells[6:7] = [20]}；每一层内部可并行，层与层之间按顺序推进。
}}
\item \code{hasCycle}: 方向是否需要有环回退路径。
\item 角方向、权重和局部方向映射。
\end{itemize}

对 230k + S128，全量 \code{orders} 大约是
\[
16640 \times 230111 \times 4 \approx 14.7\ \text{GB}
\]
仅 sweep order 就接近或超过单卡可用预算，再加场数据、工作区、metadata 会超显存。因此当前实现默认使用流式 sweep plan：一次只把一部分方向的 plan 放到 GPU。

\sect{3. SI 的 CUDA 并行方式}

\subsect{3.1 角度方向并行}

SI 的天然并行维度是角方向。给定上一轮标量通量 $\phi^{(k-1)}$，每个角方向的 sweep 彼此独立。当前实现把方向切成方向块，默认：

\[
\code{RSI\_CUDA\_SI\_PLAN\_CHUNK}=128
\]

每个方向块进入 GPU 后，执行：

\begin{enumerate}
\item 清空当前方向块的临时角通量 \code{angularPsi}。
\item 对方向块内方向做 sweep。这里 \code{previousPhi} 作为上一轮 SI 的散射源输入，不在当前方向 sweep 中被改写。
\item 每扫完一个方向块，用 \code{accumulateDirectionBatchKernel} 把该方向块的角通量按权重累加到当前轮标量通量 \code{currentPhi}。
\item 所有方向块完成后，计算相对范数并判断 SI 是否收敛。
\end{enumerate}

也就是说，SI 一轮迭代中的数据流是：

\[
\phi^{(k-1)}\ (\code{previousPhi})
\quad\longrightarrow\quad
\psi_m^{(k)}\ (\code{angularPsi})
\quad\longrightarrow\quad
\phi^{(k)}\ (\code{currentPhi})
\]

对每个 cell，累加关系为：

\[
\code{currentPhi[cell]} \mathrel{+}= w_m \cdot \code{angularPsi[m, cell]}
\]

因此扫完一个方向后，并不是更新上一轮的 \code{previousPhi}，而是把这个方向对当前轮新标量通量的贡献累加进去。所有方向贡献累加完以后，\code{currentPhi} 才成为下一轮迭代的输入。

这样避免保存完整的 $M \times C$ 角通量，只需要保存一个方向块的临时场和两个标量场缓冲区。

\subsect{3.2 无环方向的波前并行扫掠}

对无环方向，GPU 不再让一个线程串行扫完整个单元序列，而是按依赖层并行：

这里的依赖层，是按当前角方向的上风依赖关系把单元分成的一批一批集合。某个单元如果需要读取上风邻居的角通量，就必须等上风邻居先算完；没有内部上风依赖的单元放在第 0 层，只依赖第 0 层的单元放在第 1 层，依次类推。可以理解为：

\[
\text{第 0 层} \rightarrow \text{第 1 层} \rightarrow \text{第 2 层} \rightarrow \cdots
\]

同一依赖层内的单元彼此没有未满足的上风依赖，所以可以同时计算；不同依赖层之间必须按顺序推进。这就是波前并行扫掠中的波前。

\begin{itemize}
\item 代码在 CPU 端按 \code{level=0,1,2,...} 逐层 launch kernel；当前处理的是哪一层，由 launch 时传入的 \code{level} 参数决定。
\item 在普通波前 kernel \code{sweepLevelKernel} 中，\code{blockIdx.x} 对应方向块里的一个角方向。因此一个 CUDA block 处理一个角方向的当前依赖层。该层里的单元由 block 内线程按 \code{threadIdx.x} 分摊。
\item 在依赖层内单元分块 kernel \code{sweepLevelTiledKernel} 中，grid 是二维的：\code{blockIdx.x} 对应一个角方向，\code{blockIdx.y} 对应该依赖层里的一个单元分块。因此一个 CUDA block 处理一个角方向、当前依赖层、一个单元分块。
\item 这里的分块不是 \code{RSI\_CUDA\_SI\_PLAN\_CHUNK=128} 那种方向 chunk。方向 chunk 是把很多角方向分批上传和扫掠；这里的分块是把同一个依赖层里的单元再切成多个 cell tile。
\item 依赖层之间仍然保持顺序，保证上风依赖已经计算完成；依赖层内单元分块只是在同一依赖层内部增加并行 block 数。
\end{itemize}

当前主要 kernel 是 \code{sweepLevelKernel} / \code{sweepLevelTiledKernel}。分块波前并行默认开启：

\begin{itemize}
\item \code{RSI\_CUDA\_SI\_WAVEFRONT=1}
\item \code{RSI\_CUDA\_SI\_TILED\_WAVEFRONT=1}
\end{itemize}

当前每个 block 使用 \code{levelSweepThreads=256} 个线程。代码先统计当前 plan chunk 中所有方向的最大依赖层宽度 \code{maxSweepLevelWidth}，再计算
\[
\text{levelTileCount}=\left\lceil \frac{\text{maxSweepLevelWidth}}{256}\right\rceil .
\]
如果 \code{RSI\_CUDA\_SI\_TILED\_WAVEFRONT=1} 且 \code{levelTileCount>1}，SI 就使用 \code{sweepLevelTiledKernel}；否则使用普通的 \code{sweepLevelKernel}。也就是说，只有当某个依赖层宽到 256 个线程一个 block 覆盖不够时，才把依赖层切成多个单元分块。分块的目的不是改变数学算法，而是在依赖层很宽时增加 grid 维度，让更多 SM 同时工作。

\subsect{3.3 SI 中线程具体做什么}

SI 的主路径是无环波前并行扫掠。这里的 GPU 工作分配可以理解为：

\begin{itemize}
\item 不分块时，一个 CUDA block 负责一个角方向的当前依赖层；block 内线程负责该层中的不同单元。线程编号从当前依赖层的起点开始，以 \code{blockDim.x} 为步长遍历：第 0 个线程处理第 0、256、512... 个单元，第 1 个线程处理第 1、257、513... 个单元。
\item 分块时，一个 CUDA block 负责一个角方向的当前依赖层中的一个连续单元片段；通常一个线程处理该片段中的一个单元。多个 \code{blockIdx.y} 分块共同覆盖同一个依赖层。
\item 每个线程拿到一个单元后，会顺序遍历这个单元的面引用，计算该角方向下每个面的 $\mu=\Omega\cdot n$、入流项、出流项和体源项。
\item mesh 邻接关系、面几何和材料参数在进入 sweep 前已经上传到 GPU 显存。线程在 kernel 中直接从这些 global memory 数组读取，例如 \code{cellFaceOffsets}、\code{refNeighbor}、\code{refNx/refNy/refNz}、\code{refArea}、\code{volume}、\code{sigmaT}、\code{sigmaS}。
\item 如果某个入流面有内部邻居，线程用 GPU 上的 \code{refNeighbor} 找到邻居 cell id，再从当前方向的 \code{currentPsi} 中读取邻居角通量；如果是边界面，则直接调用边界入流函数。整个过程不需要回到 CPU 查询邻接关系。
\item 线程最后把该单元在当前角方向下的新角通量写入 \code{angularPsi}。
\end{itemize}

因此，在无环波前路径里，一个线程的核心任务是：计算一个角方向、一个依赖层内若干个单元的有限体积更新。线程之间不直接通信；依赖关系由依赖层顺序推进保证。

\subsect{3.4 有环方向回退路径}

如果某个方向的单元依赖图存在环，代码走 \code{sweepSamplesKernel} 回退路径。该 kernel 用投影排序扫单元，并最多做二十次局部迭代。这里的一次完整扫掠，是指按当前排序把这一方向下的所有单元遍历并更新一遍。无环方向只需要一次完整扫掠；有环方向保留最多二十次完整扫掠，用多次局部迭代缓解环形依赖。当前 230k 的高角度主路径大多是无环方向，因此主要性能来自波前并行路径。

回退路径的线程职责不同。一个 warp 有 32 个 lane，代码把它拆成 8 个四线程小组。每个四线程小组处理一个样本当前选中的方向；扫到某个四面体单元时，4 个 lane 分别计算这个单元 4 个面的入流和出流贡献。随后用 \code{warp shuffle} 在这 4 个 lane 内把贡献加起来，得到该单元的总入流和总出流，最后由第 0 个 lane 写出这个单元的新角通量。这个路径保留给有环方向或无法分层并行的情况。

\subsect{3.5 SI 的归约和收敛判断}

SI 每轮结束后做两件事：

\begin{itemize}
\item \code{accumulateDirectionBatchKernel}: 把方向通量按权重累加到 \code{currentPhi}。
\item \code{relativeNormKernel}: 在 GPU 上计算 $\|\phi^{(k)}-\phi^{(k-1)}\|_V^2$ 和 $\|\phi^{(k-1)}\|_V^2$，只拷回两个 double 判断收敛。
\end{itemize}

因此每轮迭代不会把完整场反复拷回 CPU。只有最终 SI 场会 D2H。

\sect{4. RSI 的 CUDA 并行方式}

\subsect{4.1 样本链并行}

RSI 的天然并行维度是样本。每条样本链独立选择随机方向序列。当前实现先用确定性 RNG 生成调度表：

\[
\code{seed + 17*M + 0x9e3779b9*(sample+1)}
\]

这样 CPU/GPU、样本批切分和线程调度不会改变随机链定义。

GPU 执行时按样本批分块。默认单批最多 128 个样本，同时还会根据剩余显存估算样本批容量：

\[
\text{bytes per sample} \approx 3C \cdot sizeof(double) + (N+tailExtra)\cdot sizeof(int)
\]

三个 $C$ 长度数组分别用于上一层 $\psi$、当前层 $\psi$ 和 RSI-tail 的样本局部累积。

\subsect{4.2 每批只构造实际需要的方向 plan}

S128 有 16640 个方向，但一个 RSI 样本批在 $N+10$ 步中只会用到其中一部分。当前每个样本批会：

\begin{enumerate}
\item 从调度表中收集该样本批实际用到的去重方向。
\item 对这些方向建立 global-to-local 映射。
\item 构造紧凑的 \code{DevicePlanChunk}，只上传这批方向需要的扫掠 plan。
\item 把每个迭代步/样本对应的局部方向编号上传到 \code{selectedDirections}。
\end{enumerate}

这样 RSI 不需要常驻全量 S128 plan，也不会为当前样本批没用到的方向做 GPU 扫掠。

举一个简化例子。假设当前样本批里只有 3 个样本，每条样本链只看前 4 步：

\[
\begin{array}{c|cccc}
\text{样本} & \text{第 1 步} & \text{第 2 步} & \text{第 3 步} & \text{第 4 步}\\
\hline
s_0 & 17 & 102 & 17 & 401\\
s_1 & 88 & 102 & 900 & 17\\
s_2 & 401 & 88 & 88 & 102
\end{array}
\]

这里的数字是全局角方向编号。虽然 S128 总共有 16640 个方向，但这个样本批实际只用到了
\[
\{17,88,102,401,900\}
\]
这 5 个方向。代码只为这 5 个方向准备 sweep plan，并建立局部编号，例如：
\[
17\rightarrow 0,\quad 88\rightarrow 1,\quad 102\rightarrow 2,\quad 401\rightarrow 3,\quad 900\rightarrow 4.
\]

随后调度表会被改写成局部方向编号上传到 \code{selectedDirections}。例如样本 $s_0$ 的方向序列从
\[
17,102,17,401
\]
变成
\[
0,2,0,3.
\]

GPU 推进样本 $s_0$ 时，第 1 步用方向 17 的 plan 扫一遍 mesh，得到新的 \code{currentPsi}；第 2 步再用方向 102 的 plan，以第 1 步结果作为上一层输入继续扫；第 3 步再次使用方向 17；第 4 步使用方向 401。其它样本也按自己的方向序列独立推进。这样每个样本是一条独立随机链，但同一个样本批共享这批紧凑的方向 plan。

\subsect{4.3 RSI 波前并行扫掠}

对每个 RSI 迭代步，每个样本只选一个方向。但一个方向不是只对应一个编号，而是对应这个方向下整个 mesh 的一次扫掠。对非结构网格，某个方向会把所有 cell 分成多层依赖层：
\[
\text{方向 }m \rightarrow \text{第 0 层} \rightarrow \text{第 1 层} \rightarrow \text{第 2 层} \rightarrow \cdots
\]
同一层内的 cell 可以并行，层与层之间必须按顺序推进。如果这一轮涉及的方向都无环，则 RSI 走波前并行扫掠：

\begin{itemize}
\item CPU 端仍然按 \code{level=0,1,2,...} 逐层 launch kernel。当前 launch 只处理当前依赖层。
\item 不分块时，\code{blockIdx.x} 对应样本批里的一个样本。这个 block 处理该样本当前方向的当前依赖层。block 内通常有 256 个线程，线程编号 \code{threadIdx.x} 对应该依赖层里的不同 cell：第 0 个线程处理第 0、256、512... 个 cell，第 1 个线程处理第 1、257、513... 个 cell，依次类推。
\item 分块时，grid 是二维的：\code{blockIdx.x} 仍然对应一个样本，\code{blockIdx.y} 对应该依赖层里的一个 cell tile。此时一个 block 只处理当前依赖层中的一段连续 cell，通常一个线程处理这个 tile 里的一个 cell；多个 \code{blockIdx.y} 共同覆盖同一个样本、同一个方向、同一个依赖层。
\item block 内线程之间不需要交换角通量。每个线程写自己负责的 cell；层间依赖由 CPU 端按 \code{level} 顺序 launch kernel 来保证。
\item 如果两个样本同一步选到同一个方向，它们可以共享这个方向的 sweep plan，但 \code{previousPsi/currentPsi} 按样本分开存储，所以不会互相覆盖。
\end{itemize}

RSI 的依赖层内单元分块还多一个样本批条件。只有同时满足以下条件时才分块：
\begin{itemize}
\item \code{RSI\_CUDA\_RSI\_TILED\_WAVEFRONT=1}；
\item \code{levelTileCount>1}，即最大依赖层宽度超过 256 个单元；
\item \code{batchSize>=64}，即样本批本身足够大。
\end{itemize}
否则普通波前并行已经能提供足够的 grid 规模。有环方向单独处理：
\begin{itemize}
\item 如果该迭代步中有有环方向，则该迭代步回退到 \code{sweepSamplesKernel}。
\item 这是保守策略，保证依赖处理简单。
\end{itemize}

\subsect{4.4 RSI 与 RSI-tail 合并计算}

当前实现不会分别跑 RSI 和 RSI-tail。对每个样本：

\begin{itemize}
\item 第 $N$ 步的 \code{currentPsi} 进入 \code{globalRSISum}。
\item 第 $N$ 到 $N+10$ 步的 \code{currentPsi} 累加到 \code{sampleTail}。
\item 一个样本批结束后，再把 \code{sampleTail} reduce 到 \code{globalTailSum}。
\end{itemize}

这样 tail 的 11 个时间层只在样本局部工作区中累加，避免每个 tail 迭代步都做一次全局 cell reduction。最终：

\[
\phi_{\mathrm{RSI}} = \frac{1}{S}\sum_s \psi_s^{(N)},\qquad
\phi_{\mathrm{tail}} = \frac{1}{S(10+1)}\sum_s \sum_{j=0}^{10}\psi_s^{(N+j)}
\]

\sect{5. 显存与流式 Plan 策略}

当前入口会估算全量 \code{orders} 的显存：

\[
M \times C \times sizeof(int)
\]

如果超过阈值，默认 8 GiB，就走流式 sweep plan。该阈值可用：

\[
\code{RSI\_CUDA\_STREAMING\_PLAN\_THRESHOLD\_MB}
\]

调整。

流式 sweep plan 的核心逻辑是：

\begin{itemize}
\item mesh、ordinates、weights 等静态数据只上传一次。
\item SI 按固定方向块顺序处理。
\item RSI 按样本批实际用到的去重方向临时组装紧凑 plan。
\item plan 不直接等同于场数据；场数据工作区与 plan 工作区分开管理。
\end{itemize}

这使得 230k + S128 + 8192 样本可以在 16 GB 级别 GPU 上完整运行。

\sect{6. 缓存与 Plan 优化逻辑}

\subsect{6.1 磁盘方向块缓存}

每个角方向在 sweep 前都要先生成一个 sweep plan。它记录这个方向下 GPU 扫掠需要的结构信息，包括：

\begin{itemize}
\item cell 应该按什么顺序扫；
\item 该方向的依赖图是否有环；
\item 无环方向有哪些依赖层；
\item 每个依赖层在 sweep order 中从哪里开始、到哪里结束。
\end{itemize}

生成这些 plan 需要遍历 mesh 的邻接关系。对 230k + S128 来说，方向数达到 16640，如果每次运行都重新构建所有方向的 plan，开销会很大。因此当前实现把已经生成好的方向块 plan 序列化到磁盘：

\[
\code{results/cache/rsi\_sweep\_plan\_chunk\_<key>.bin}
\]

下次遇到同样的 mesh、同样的角方向集合和同样的缓存格式版本时，就直接从这个 \code{.bin} 文件读回 plan，不重新构建。cache key 由 mesh、角方向、方向列表和格式版本决定；如果 mesh 变了、方向集合变了，或者缓存文件损坏/版本不匹配，代码会忽略旧缓存并重新构建。

SI 和 RSI 的缓存粒度不同。SI 每轮本来就要顺序扫完整的方向集合，因此用较大的方向块比较合适；当前默认 \code{RSI\_CUDA\_SI\_PLAN\_CHUNK=128}，S128 的 16640 个方向大约对应 130 个 SI 方向块缓存。RSI 则不同，每个样本批只随机用到其中一部分方向。例如 128 个样本的 batch 往往只用到约 2800 个去重方向，而且这些方向随机分布在 0 到 16639 之间。如果 RSI 也按 128 个方向一块缓存，一个 batch 只要碰到某个 chunk 中的少数方向，就可能需要读取整个 128-direction chunk，最后接近读取全量 S128 plan，包含大量本批没用到的方向。

因此当前 RSI 固定方向缓存默认使用
\[
\code{RSI\_CUDA\_FIXED\_PLAN\_CHUNK=1}
\]
也就是一个缓存文件对应一个方向。这样 S128 最多约 16640 个方向缓存文件，但每个样本批只加载自己实际用到的方向，避免随机访问时读入大量无用 plan。

\subsect{6.2 主机端 LRU 缓存}

磁盘缓存解决的是 ``不要重新生成 plan''，但它没有消除读文件本身的开销。如果某个方向块每次都从 \code{.bin} 文件取，仍然要经历：

\begin{enumerate}
\item 打开缓存文件；
\item 读取二进制内容；
\item 检查 header、方向数和 cell 数；
\item 反序列化出每个方向的 \code{order}、\code{levelOffsets} 和 \code{hasCycle}；
\item 再把这些 \code{SweepPlan} 打包成 GPU 上传需要的 \code{HostPlanChunk}。
\end{enumerate}

所以当前实现又加了一层 CPU 内存缓存，也就是主机端 LRU 缓存：

\[
\code{RSI\_CUDA\_PLAN\_HOST\_CACHE\_MB}
\]

默认 1024 MiB。这里缓存的不是 GPU 设备指针，而是主机内存里的 \code{vector<SweepPlan>}。缓存键和磁盘 \code{.bin} 文件使用同一个方向块 key；也就是说，同一个 mesh、同一个方向列表、同一个缓存格式版本，会查到同一个主机端缓存项。

一次方向块准备的查找顺序是：

\begin{enumerate}
\item 先用 cache key 查主机端 LRU。
\item 如果命中，直接复制内存中的 \code{vector<SweepPlan>}，跳过磁盘读取和反序列化。
\item 如果没有命中，再去读对应的 \code{.bin} 文件。
\item 如果磁盘也没有命中，才重新遍历 mesh 邻接关系构建 sweep plan。
\item 从磁盘读到或者重新构建出的 plan，会再放回主机端 LRU，供后续方向块复用。
\end{enumerate}

LRU 是 least recently used，中文可以理解成 ``最近最少使用''。当缓存占用超过预算时，代码优先丢掉最久没有被访问的方向块，保留最近反复用到的方向块。这样做的原因是 SI 和 RSI 都有明显的重复访问：

\begin{itemize}
\item SI 每轮迭代都按相同方向块顺序扫完整角方向集合。第 1 轮把方向块放入 LRU 后，后面的 SI 轮次再次访问同一个 chunk 时，通常可以直接从主机内存拿到 plan。
\item RSI 每个样本批只使用一部分随机方向，但不同样本批之间会反复碰到相同的固定方向。单方向固定缓存配合主机端 LRU 后，后续 batch 再需要这个方向时，也可以避免重复读磁盘。
\end{itemize}

需要注意的是，主机端 LRU 只减少 CPU 侧准备 plan 的开销；它本身不等于 GPU 端常驻。命中 LRU 后，代码仍然要把这些 plan 重新组装成 \code{HostPlanChunk}，并在需要时上传成 \code{DevicePlanChunk}。如果想进一步避免 SI 每轮重复 H2D 上传，就要靠下一节的 SI 设备端 plan 缓存。

因此缓存层级可以理解为：

\[
\text{重新构建 plan} \;>\; \text{读磁盘 .bin} \;>\; \text{主机端 LRU 命中} \;>\; \text{GPU 设备端常驻}
\]

越往右复用成本越低。

\subsect{6.3 SI 设备端 plan 缓存}

SI 每轮迭代都会按同一组角方向、同一个方向块顺序重复扫掠。例如 S128 使用 \code{RSI\_CUDA\_SI\_PLAN\_CHUNK=128} 时，16640 个方向会被切成约 130 个方向块。第 1 轮 SI 要扫这 130 个方向块；第 2 轮、第 3 轮也还是按同样顺序再扫同一批方向块。

如果只有 6.2 的主机端 LRU，那么第 2 轮以后可以少读磁盘，但仍然要把每个方向块从 CPU 侧组装并上传到 GPU。对 SI 来说，这个上传会在每一轮重复发生。因此当前实现又提供了一层 SI 专用的设备端 plan 缓存：

\[
\code{RSI\_CUDA\_SI\_DEVICE\_PLAN\_CACHE\_MB}
\]

默认 8192 MiB。这里缓存的是已经上传到 GPU 的 \code{DevicePlanChunk}，包括该方向块在 GPU sweep 中要用的方向编号、方向余弦、权重、cell 扫掠顺序、依赖层 offset 等数组。

一次 SI 方向块的处理流程是：

\begin{enumerate}
\item CPU 按 \code{chunkIndex} 访问当前方向块。
\item 先检查 \code{siDevicePlanCache[chunkIndex]} 是否已经有对应的 \code{DevicePlanChunk}。
\item 如果命中，直接用 GPU 上已有的 plan launch sweep kernel，不需要重新上传这个方向块的 plan。
\item 如果没有命中，就走 6.2 的主机端准备流程：查 LRU、查磁盘，必要时重新构建 \code{SweepPlan}。
\item 主机端 plan 准备好后，上传成 \code{DevicePlanChunk}。
\item 如果该方向块不超过设备端缓存预算，并且剩余预算够用，就把它留在 \code{siDevicePlanCache} 对应位置；否则只临时使用，用完释放。
\end{enumerate}

因此第 1 轮 SI 通常会产生较多 cache miss，因为方向块第一次出现时还没在 GPU 上。只要显存预算够，代码会把这些上传过的方向块留下。第 2 轮以后再次访问相同 \code{chunkIndex} 时，就变成 cache hit，直接复用 GPU 内存里的 plan。

这个策略主要适合 SI，而不是 RSI，原因是 SI 的方向访问顺序固定并且重复很多轮；RSI 每个样本批的随机方向集合不同，设备端常驻大块 plan 的命中规律没有 SI 稳定。RSI 主要依赖 6.4 的固定方向缓存和 6.5 的预热/并行加载来降低 plan 准备开销。

简单说，三层缓存各自解决的问题不同：

\begin{itemize}
\item 磁盘缓存：避免重新遍历 mesh 构建 sweep plan。
\item 主机端 LRU：避免反复读 \code{.bin} 和反序列化。
\item SI 设备端缓存：避免 SI 每轮把同一个方向块 plan 反复从 CPU 上传到 GPU。
\end{itemize}

\subsect{6.4 RSI 固定方向缓存}

RSI 的方向访问模式和 SI 不一样。SI 是固定顺序反复扫完整角方向集合；RSI 是每个样本链按调度表随机选方向。以 S128、8192 个样本为例，总方向数是 16640，但每个 128-sample batch 通常只用到约 2800 个去重方向，而且这些方向在 0 到 16639 之间随机分布。

如果 RSI 直接按 128 个方向一个 chunk 来复用 plan，会有一个问题：某个 batch 可能只需要 chunk 里的少数几个方向，却必须把整个 128-direction chunk 都读出来、打包、上传。因为随机方向覆盖范围很散，一个 batch 最后可能触碰很多 chunk，等价于加载大量本批不用的方向。

所以当前 RSI 使用固定方向缓存，把大方向块拆成更小的固定方向块。默认：

\[
\code{RSI\_CUDA\_FIXED\_PLAN\_CHUNK}=1
\]

也就是一个固定缓存块只包含一个全局角方向。对 S128 来说，最多约 16640 个固定方向缓存文件。这样文件数量变多，但每个样本批只加载自己真正用到的方向，避免随机访问时把大量无关方向带进来。

一个 RSI batch 准备 plan 时，大致流程是：

\begin{enumerate}
\item 从该 batch 的所有样本链和所有迭代步中收集全局方向编号。
\item 排序并去重，得到 \code{uniqueDirections}。
\item 对每个需要的全局方向，找到它所属的固定方向缓存块。
\item 对已经有缓存的方向块，直接读固定方向 plan。
\item 对没有缓存的方向块，重新构建对应方向的 \code{SweepPlan}，然后保存到磁盘缓存。
\item 把这些固定方向 plan 按当前 batch 的 \code{uniqueDirections} 顺序重新组装，形成紧凑的 \code{HostPlanChunk}。
\item 再上传成当前 batch 使用的 \code{DevicePlanChunk}。
\end{enumerate}

这里要注意，固定方向缓存解决的是 ``方向 plan 的复用''，不是让每个方向都常驻 GPU。RSI 每个 batch 仍然会构造一个只包含本批去重方向的紧凑 \code{DevicePlanChunk}。这样 GPU sweep 时只面对当前 batch 实际需要的方向，内存占用比全量 S128 plan 小得多。

日志里的 \code{rsi\_plan\_fixed\_cache\_hits} 和 \code{rsi\_plan\_fixed\_cache\_misses} 就是在看这层固定方向缓存的效果。hit 表示这个固定方向 plan 已经存在，可以复用；miss 表示这个方向还没有缓存，需要构建并保存。

\subsect{6.5 RSI 预热与并行加载}

固定方向缓存能减少重复构建，但如果等到每个 batch 真正运行时才发现某个方向 miss，就会把构建 plan 的开销插入到 batch 准备阶段。为了减少运行中突然 miss，当前 RSI 在正式 batch sweep 前可以做预热：

\[
\code{RSI\_CUDA\_RSI\_PLAN\_PREWARM=1}
\]

默认开启。预热的前提是 RSI 方向调度表已经生成。调度表生成后，代码已经知道所有样本、所有迭代步会访问哪些全局方向，因此可以在正式 sweep 前先扫描完整调度表。这里得到的不是某一个 batch 的 \code{uniqueDirections}，而是整次 run 会用到的全局方向集合，再把这些方向映射成固定方向块集合。

如果 \code{RSI\_CUDA\_FIXED\_PLAN\_CHUNK=1}，一个固定方向块正好就是一个全局方向，所以 ``固定方向块集合'' 和 ``整次 run 的去重方向集合'' 基本等价。如果固定方向块大小大于 1，一个固定方向块会包含多个连续全局方向；只要其中任意一个方向被本次 run 用到，这个固定方向块就会进入预热集合。

预热阶段的流程是：

\begin{enumerate}
\item 扫描完整 RSI 调度表，先得到整次 run 的去重全局方向集合。
\item 把这些全局方向映射到固定方向块，得到需要预热的固定方向块集合。
\item 对每个需要预热的固定方向块，先检查对应的磁盘缓存文件是否已经存在。
\item 如果文件已经存在，只记录为 prewarm hit，不重新构建。
\item 如果文件不存在，记录为 prewarm miss，并构建这个固定方向的 sweep plan。
\item 构建完成后保存为固定方向缓存文件，后续 batch 直接读取。
\end{enumerate}

这一步的目的不是提前把所有 plan 放进 GPU，而是提前补齐磁盘/主机侧可复用的固定方向 plan。这样 batch 阶段更可能只做读取、组装和上传，而不是边运行边构建。

预热和固定方向加载可以并行执行：

\[
\code{RSI\_CUDA\_FIXED\_PLAN\_PARALLEL\_LOAD=1}
\]

默认开启。实际含义是用多个 CPU worker 同时处理多个固定方向块：有缓存的方向并行读取；缺缓存的方向并行构建并保存。预热 worker 数可以由 \code{RSI\_CUDA\_FIXED\_PLAN\_PREWARM\_WORKERS} 控制；如果没有单独设置，通常复用固定方向加载 worker 的设置。

还需要 miss 去重。因为多个 batch 或多个 worker 可能同时发现同一个固定方向缺缓存，如果不去重，就可能有多个线程重复构建同一个方向 plan。当前逻辑默认开启 miss 去重：

\[
\code{RSI\_CUDA\_FIXED\_PLAN\_MISS\_DEDUP=1}
\]

含义是：同一个固定方向块 miss 时，只让一个 worker 负责构建和保存，其它 worker 等待；等缓存文件生成后，其它 worker 再读取结果。这样减少重复构建，也避免多个线程同时写同一个缓存文件。

因此 6.4 和 6.5 配合起来看，RSI plan 优化分成两步：

\begin{itemize}
\item 固定方向缓存：把随机方向访问拆成可复用的单方向 plan，避免读入无关方向。
\item 预热与并行加载：在正式 batch sweep 前尽量补齐这些单方向 plan，并把缺失构建并行化。
\end{itemize}

以现有 230k + S128 + 8192 日志为例，预热阶段输出：

\[
\code{CUDA RSI fixed plan prewarm: chunks=16640, hits=16639, misses=1}
\]

说明这次运行前 S128 的固定方向缓存基本已经齐全，正式 RSI batch 中主要成本就转移到固定方向 plan 的读取、组装、上传和真正 sweep 上，而不是大量重新构建 sweep plan。

\subsect{6.6 RSI fused wavefront sweep}

6.4 和 6.5 主要降低 RSI 的 plan 准备成本；当固定方向缓存已经齐全以后，230k + S128 + 8192 的主开销又回到真正的 sweep kernel。当前 RSI 对无环方向提供 fused wavefront 路径：

\[
\code{RSI\_CUDA\_RSI\_FUSED\_WAVEFRONT}=1
\]

每个 RSI batch 包含若干条样本链，每个样本在当前 iteration 只选择一个方向。对无环方向，当前 fused 路径把该样本当前方向的所有依赖层放在同一个 kernel 内部顺序推进，避免在 CPU 端为每一层单独组织一次 launch。230k 网格的每个方向通常有数百个依赖层；对 8192 样本、$N+10$ 个随机步来说，level 组织方式会直接影响 sweep 阶段的调度成本和尾部并行度。

该路径复用已有的 \code{sweepLevelFusedKernel}。一个 CUDA block 仍然对应一个样本；block 内线程处理该样本当前方向的一个依赖层。不同的是，kernel 内部按 level 顺序循环：

\[
\code{level=0} \rightarrow \code{level=1} \rightarrow \cdots
\]

每层结束后用 \code{\_\_syncthreads()} 保证同一个样本、同一个方向内部的上风依赖已经完成，然后继续下一层。这样一次 kernel launch 就覆盖一个样本当前方向的一整次无环 sweep，而不是每个 level 单独 launch。

这个 fused 路径只用于无环 RSI wavefront；如果某个 iteration 中存在有环方向，仍然回退到 \code{sweepSamplesKernel} 多 pass 路径。当前默认策略是：

\begin{itemize}
\item 当 cell 数 $C \ge 100000$ 时默认启用 fused RSI wavefront，因为大网格的 level 数多，减少 launch 数更有收益。
\item 小网格默认保留 tiled level-by-level 路径，避免在 level 数较少时引入不必要的 fused kernel 组织成本。
\item 可用 \code{RSI\_CUDA\_RSI\_FUSED\_WAVEFRONT=0} 显式关闭，便于定位性能或数值问题。
\end{itemize}

从工作分配看，fused RSI wavefront 没有改变随机链、方向调度表、sweep plan、边界源或离散方程。它只改变同一个样本当前方向内部 level 的 launch 组织方式：

\begin{itemize}
\item level-by-level 路径：CPU 端逐 level launch，level 间顺序由 kernel launch 顺序保证。
\item fused 路径：GPU kernel 内逐 level 循环，level 间顺序由同一 block 内的 \code{\_\_syncthreads()} 保证。
\end{itemize}

因此 fused 路径和 level-by-level 路径应当给出逐值一致的结果。实际验证中，230k + S64 + 512 的 \code{SI\_fine.csv}、\code{RSI.csv} 和 \code{RSI\_tail.csv} 在 fused 开关前后逐字节一致；小规模 \code{gpu\_consistency\_test} 的 CPU/GPU 相对 $L_2$ 误差约为 $10^{-16}$。

在当前 230k + S128 + 8192 streaming 日志中，RSI fused wavefront 运行在 S128 RSI sweep 阶段，当前口径为：

\begin{tabularx}{\linewidth}{>{\bfseries}l r r r}
\toprule
指标 & total / s & plan / s & sweep / s\\
\midrule
\rowcolor{softrow}
RSI & 701.386 & 97.775 & 566.236\\
\bottomrule
\end{tabularx}

这说明当前完整 Figure 5 时间中，RSI 仍主要由 sweep 本身和 plan 组装/上传组成。fused wavefront 覆盖的是无环 RSI sweep 的 level 组织方式；它不处理 SI sweep，也不消除 RSI plan 的固定方向读取、assemble、pack 和 upload 成本。

这次也尝试过把 \code{sigmaT * volume}、\code{sigmaS * volume} 和 \code{q * volume} 预计算到 GPU mesh 数组中，减少每个 cell 更新里的乘法。这里的 A/B 指两个实现变体：A 为当前保留路径，在 sweep kernel 内按需计算这些乘积；B 为实验路径，在 mesh 上传阶段预先生成乘积数组，sweep kernel 内直接读取。230k + S16 + 16 的快速验证结果波动较大，fine SI sweep 没有稳定改善，甚至有小幅回归。因此该方案没有保留。当前判断是 sweep kernel 更受 wavefront 并行度、global memory 访问和 kernel launch 组织影响，单元公式中的少量算术不是主要瓶颈。

\sect{7. Mesh 数据布局优化}

这一节的核心是：把 sweep 计算 cell 时需要的 mesh 数据，提前整理成 GPU 友好的数组，并一次性上传到 GPU global memory。线程处理 cell 时直接读取 \code{volume}、\code{sigmaT}、\code{sigmaS}、\code{cellQ}、\code{cellFaceOffsets}、\code{refNeighbor}、\code{refNx/refNy/refNz}、\code{refArea} 和边界数组，不需要运行中回到 CPU 查询 mesh。

当前 GPU mesh 布局类似 CSR：\code{cellFaceOffsets[cell]} 到 \code{cellFaceOffsets[cell+1]} 给出该 cell 的面引用区间；每个 \code{refIndex} 在多条 SoA 数组中保存邻居、法向、面积、面中心和边界信息。这样 kernel 不必从 cell 再追到 face 对象，减少了一层间接索引，也降低了 GPU 端对象结构开销。

需要注意的是，这并不把非结构网格访存变成规则连续访问。不同线程处理的 cell 在 sweep order 中不一定相邻，读取 \code{currentPsi[neighbor]} 仍可能是随机位置。这个优化的作用是把几何、邻接和边界数据放在 GPU 上并按面引用对齐，让 SI、RSI 和 RSI-tail 共用同一份 GPU mesh 数据。

\sect{8. 当前性能口径}

当前 Figure 5 日志包含两组网格、两组角度口径。S4 是 coarse SI，$M=24$，只做 SI-only；S128 是 fine SI + RSI + RSI-tail，$M=16640$，RSI 使用 8192 条样本链。下面按 30k non-streaming、30k streaming 和 230k streaming 三种当前运行路径列出日志口径。

\subsect{8.0 计时字段说明}

\begin{tabularx}{\linewidth}{>{\bfseries}l X}
\toprule
字段 & 含义\\
\midrule
\code{key} & 计算 sweep plan 缓存 key 和缓存文件路径的时间。\\
\code{cache} & 从已有 plan cache 读取并反序列化的时间；RSI fixed cache 中这是多线程累计时间，不一定等于 wall time。\\
\code{build} & cache miss 时从 mesh 和角方向重新构建 \code{SweepPlan} 的时间，包括计算 \code{order}、\code{levelOffsets} 和 \code{hasCycle}。\\
\code{save} & 把新构建的 sweep plan 写入磁盘 cache 的时间。\\
\code{assemble} & 把固定方向 plan 按当前 batch 的 \code{uniqueDirections} 顺序组装成紧凑 plan 的时间，主要用于 RSI。\\
\code{pack} & 把 \code{SweepPlan} 打包成可上传的 \code{HostPlanChunk} 数组布局的时间。\\
\code{upload} & 把 mesh、方向、plan 或工作数组上传到 GPU 的时间，具体含义随日志阶段而定。\\
\code{sweep} & GPU 上真正执行输运 sweep kernel 的时间，即按方向/样本和依赖层计算 cell 通量。\\
\code{schedule} & 生成 RSI 随机方向调度表的时间。\\
\code{unique} & 每个 RSI sample batch 从调度表中收集、排序并去重实际用到方向的时间。\\
\code{reduce}/\code{norm}/\code{copy} & 分别是结果归约、SI 收敛范数计算和最终 GPU 到 CPU 拷贝时间。\\
\bottomrule
\end{tabularx}

\subsect{8.1 30k non-streaming 计时}

\begin{tabularx}{\linewidth}{>{\bfseries}l r r r r r r}
\toprule
角度 & 方向 & plan/setup / s & SI total / s & SI sweep / s & RSI total / s & RSI sweep / s\\
\midrule
\rowcolor{softrow}
S4 & 24 & 0.160 & 0.219 & 0.049 & -- & --\\
S128 non-streaming & 16640 & 1.738 & 60.139 & 58.034 & 88.500 & 84.712\\
\bottomrule
\end{tabularx}

这里的 \code{plan/setup} 采用 CUDA 日志中的 upload 口径。non-streaming 的 S128 全量 plan 可以放入显存，fine SI 和 RSI/RSI-tail 共享同一套 GPU 问题数据；日志不会把这部分拆成独立的 \code{SI plan} 和 \code{RSI plan}。

为了测量从零准备完整 sweep plan 的时间，另用临时缓存目录
\[
\code{RSI\_SWEEP\_PLAN\_CACHE\_DIR=/tmp/rsi\_sweep\_cache\_cold\_30k\_nonstream}
\]
重新跑了一次 30k non-streaming。此时原有 \code{results/cache} 不参与命中，日志中的完整 plan 准备时间为：

\begin{tabularx}{\linewidth}{>{\bfseries}l r r r r}
\toprule
阶段 & 方向 & build / s & save / s & 说明\\
\midrule
\rowcolor{softrow}
coarse SI plan & 24 & 0.010 & 0.003 & 构建并缓存 S4 sweep plan。\\
fine SI/RSI plan & 16640 & 9.203 & 2.742 & 构建并缓存 S128 全量 sweep plan。\\
\bottomrule
\end{tabularx}

同一次冷缓存运行中，S128 GPU problem 上传 \code{upload=1.847 s}，RSI 随机方向调度表生成 \code{rsi\_schedule=0.485 s}，外部 wall time 为 \code{elapsed 159.77 s}。因此对 30k non-streaming 来说，一开始完整计算扫掠顺序和依赖层的主要成本就是上表的 \code{build}，而不是后续 batch 中的 fixed-direction miss。

30k non-streaming 这次完整命令的外部 wall time 为：

\[
\code{elapsed\ 151.05\ s}
\]

\subsect{8.2 30k streaming 计时}

为了让 30k 也走和 230k 相同的 streaming 逻辑，强制设置
\[
\code{RSI\_CUDA\_STREAMING\_PLAN\_THRESHOLD\_MB=0}
\]
跑了一次 30k S128 streaming。这个路径会像 230k 一样把 SI 的 16640 个方向组织成 130 个 128-direction chunk，并让 RSI 每个 batch 构造紧凑 plan。

\begin{tabularx}{\linewidth}{>{\bfseries}l r r r r r r}
\toprule
角度 & 方向 & plan/setup / s & SI total / s & SI sweep / s & RSI total / s & RSI sweep / s\\
\midrule
\rowcolor{softrow}
S4 & 24 & 0.148 & 0.204 & 0.046 & -- & --\\
S128 streaming & 16640 & 19.299 & 60.965 & 52.591 & 102.494 & 84.747\\
\bottomrule
\end{tabularx}

30k streaming 的 plan/setup 可以拆成：

\begin{tabularx}{\linewidth}{>{\bfseries}l r r r r r r}
\toprule
阶段 & total / s & key / s & cache / s & assemble / s & pack / s & upload / s\\
\midrule
\rowcolor{softrow}
SI plan & 6.356 & 0.008 & 5.169 & 0.000 & 0.156 & 0.880\\
RSI plan & 12.943 & 1.137 & 130.667 & 3.841 & 1.734 & 8.289\\
\bottomrule
\end{tabularx}

这里 30k streaming 的 \code{plan/setup} 为 19.299 s，由 6.356 s 的 SI plan 和 12.943 s 的 RSI plan 组成。其中 RSI plan 的 cache 计时为 130.667 s，这是多线程固定方向缓存读取的累计计时，不等于 wall time；日志中对应的固定方向缓存 wall time 是 8.577 s。

30k streaming 这次完整命令的外部 wall time 为：

\[
\code{elapsed\ 158.44\ s}
\]

30k streaming 中 SI 的 130 个 plan chunk 全部常驻设备端缓存：

\[
\code{device\ chunks=130/130,\ hits=1690,\ misses=130}
\]

因此 30k streaming 的 SI sweep 主要受 GPU sweep kernel 本身影响；RSI 仍需要按样本批读取固定方向 plan、组装紧凑 plan 并上传。

\subsect{8.3 230k streaming 计时}

\begin{tabularx}{\linewidth}{>{\bfseries}l r r r r r r}
\toprule
角度 & 方向 & plan/setup / s & SI total / s & SI sweep / s & RSI total / s & RSI sweep / s\\
\midrule
\rowcolor{softrow}
S4 & 24 & 0.305 & 0.565 & 0.228 & -- & --\\
S128 streaming & 16640 & 145.147 & 801.394 & 720.885 & 701.386 & 566.236\\
\bottomrule
\end{tabularx}

230k 的 S128 走 streaming plan 路径，因此日志可以单独统计 SI plan 和 RSI plan。本机约 16 GB host memory 下，S128 + 8192 使用

\[
\code{RSI\_CUDA\_RSI\_MAX\_SAMPLES\_PER\_BATCH=64}
\]

限制 RSI batch 大小，避免 host 侧 batch plan/prefetch 内存峰值过高。

\begin{tabularx}{\linewidth}{>{\bfseries}l r r r r r r}
\toprule
阶段 & total / s & key / s & cache / s & assemble / s & pack / s & upload / s\\
\midrule
\rowcolor{softrow}
SI plan & 47.372 & 0.046 & 16.079 & 0.000 & 1.480 & 24.122\\
RSI plan & 97.775 & 3.557 & 955.513 & 41.186 & 10.759 & 53.701\\
\bottomrule
\end{tabularx}

注意这里的 \code{RSI plan cache=955.513 s} 是多次固定方向缓存读取的累计计时，不等于 wall time；日志中对应的固定方向缓存 wall time 是 64.368 s。230k streaming 中 SI 设备端缓存不能容纳全部 130 个 chunk，本次日志为：

\[
\code{device\ chunks=88/130,\ host\ chunks=36/130}
\]

230k streaming 这次完整命令的外部 wall time 为：

\[
\code{elapsed\ 1437.14\ s}
\]

综合这三个小节可以看到：S4 coarse SI 很快；S128 下真正占主导的是 sweep。30k non-streaming 的 S128 中，SI sweep 约 58.0 s，RSI sweep 约 84.7 s；30k streaming 的 S128 中，SI sweep 约 52.6 s，RSI sweep 约 84.7 s；230k streaming 的 S128 中，SI sweep 约 720.9 s，RSI sweep 约 566.2 s。tail 累加、最终 reduce、方向调度和 D2H copy 都不是主瓶颈。plan 相关成本已经通过缓存、预热和 streaming plan 控制在可运行范围内，但 230k + S128 的 plan 读取/组装/上传仍然是可见开销。

\subsect{附录 A：历史 27 组 RSI/SI 并行效率测试}

本附录保留早期验证数据；当前完整因子矩阵及同网格参考误差定义见后面的 8.4。

以下两表覆盖全部 27 个已完成的空间--角度--样本--batch 组合。新增 S100（10,200 个方向）在 10k/30k/100k 三种网格上以 5,000 样本、$B=128$ 复测，覆盖约 $10^2$、$10^3$、$10^4$ 三个角度规模；10k/30k 的 S10--S30 历史日志也已纳入。100k/S224 的 100--100,000 样本与 $B=32/128/256$ 扫描保留，用于分离样本数和 batch 的影响。网格标签后的数值是实际 cell 数量的约数：10k/9k、30k/36k、100k/94k。$B$ 为实际 batch capacity；`--' 表示旧日志未记录该字段。

表中 error 为 RSI-tail 最终场相对同一空间网格、同一 rep=1 的确定性 SI 最终场的相对 $L_2$ 误差。每个参考场均采用 S316，即 100,488 个方向；参考 SI 的具体配置如下。SI 在 RSI 前独立完成，故下表的关联样本数和 $B$ 只是同一日志中随后 RSI 的设置，不会影响参考 SI 场。

{\scriptsize\setlength{\tabcolsep}{4pt}
\begin{tabular}{lrrr}
\toprule
参考空间网格/cells & 参考 SI 角度/方向数 & 关联样本数 & $B$\\
\midrule
10k/9k & S316 / 100488 & 5000 & 128\\
30k/36k & S316 / 100488 & 5000 & 128\\
100k/94k & S316 / 100488 & 5000 & --\\
\bottomrule
\end{tabular}}

{\scriptsize\setlength{\tabcolsep}{3pt}
\begin{tabular}{lrrrrrrr}
\toprule
网格/cells & 角度/方向数 & 样本数 & $B$ & SI total/s & RSI total/s & RSI/SI & error\\
\midrule
10k/9k&120&5000&128&0.042&2.276&0.02x&12.556\%\\
10k/9k&440&5000&128&0.149&2.831&0.05x&6.497\%\\
10k/9k&960&5000&128&0.285&3.350&0.08x&4.533\%\\
10k/9k&5040&1000&--&1.634&1.770&0.92x&3.882\%\\
10k/9k&5040&5000&128&1.094&4.845&0.23x&1.920\%\\
10k/9k&10200&5000&128&2.491&5.763&0.43x&1.796\%\\
10k/9k&50624&5000&128&14.572&14.031&1.04x&0.956\%\\
10k/9k&100488&5000&128&29.554&18.841&1.57x&0.842\%\\
30k/36k&120&5000&128&0.143&5.933&0.02x&11.211\%\\
30k/36k&440&5000&128&0.457&6.897&0.07x&6.037\%\\
30k/36k&960&5000&128&18.495&256.681&0.07x&4.132\%\\
30k/36k&10200&5000&128&44.393&65.697&0.68x&1.680\%\\
30k/36k&50624&5000&128&169.717&63.580&2.67x&0.888\%\\
30k/36k&100488&5000&128&399.769&75.389&5.30x&0.757\%\\
100k/94k&10200&5000&128&14.666&40.123&0.37x&1.892\%\\
100k/94k&50624&100&100&136.412&1.694&80.51x&14.289\%\\
100k/94k&50624&500&32&133.890&6.225&21.51x&5.475\%\\
100k/94k&50624&500&128&136.112&6.263&21.73x&5.475\%\\
100k/94k&50624&500&256&142.403&8.213&17.34x&5.475\%\\
100k/94k&50624&1000&128&136.483&11.034&12.37x&2.697\%\\
100k/94k&50624&5000&--&221.185&133.404&1.66x&1.062\%\\
100k/94k&50624&5000&128&138.568&49.903&2.78x&1.062\%\\
100k/94k&50624&10000&128&135.631&117.637&1.15x&0.782\%\\
100k/94k&50624&50000&128&141.423&499.754&0.28x&0.344\%\\
100k/94k&50624&100000&128&135.101&980.555&0.14x&0.315\%\\
100k/94k&100488&5000&--&719.068&246.544&2.92x&0.902\%\\
100k/94k&100488&10000&128&639.425&168.518&3.79x&0.682\%\\
\bottomrule
\end{tabular}}

表 1 显示，低角度下 RSI 的 batch、随机调度和 plan 开销不能摊销；120 个方向时三种网格均慢于 SI。高角度组合则出现总时间收益，例如 30k/S316/5,000 样本为 5.30x（0.757\%），100k/S316/10,000 样本为 3.79x（0.682\%）。S100 的总时间均使用 fixed-plan cache 预热后的复测结果，避免把首次建 cache 的成本混入网格间对比。

两种方法按同一个 directional cell-sweep 工作单位计算并行效率，单位均为 Mcell-sweeps/s：
\[
E_{\rm SI}=\frac{\mathrm{cells}\times N_{\Omega}\times N}{t_{\rm SI,sweep}\times10^6},\qquad
E_{\rm RSI}=\frac{\mathrm{cells}\times N_s\times(N+10)}{t_{\rm RSI,sweep}\times10^6}.
\]
其中 $N_{\Omega}$ 为角度/方向数，$N_s$ 为 RSI 样本数，$N$ 是 SI 收敛迭代数，$N+10$ 对应 RSI-tail 共用链的步数。$E_{\rm RSI}/E_{\rm SI}$ 是在同一工作单位下 RSI 相对 SI 的 GPU 并行效率提升，而不是端到端时间加速比；plan/s 保留以识别准备和传输成本。

{\scriptsize\setlength{\tabcolsep}{2pt}
\begin{tabular}{lrrrrrrrrr}
\toprule
网格/cells&角度/方向数&样本数&$B$&SI sweep/s&$E_{\rm SI}$&RSI sweep/s&$E_{\rm RSI}$&$E_{\rm RSI}/E_{\rm SI}$&plan/s\\
\midrule
10k/9k&120&5000&128&0.011&1449.98&1.892&582.65&0.40x&0.110\\
10k/9k&440&5000&128&0.041&1380.92&2.066&533.62&0.39x&0.459\\
10k/9k&960&5000&128&0.091&1353.64&2.107&523.20&0.39x&0.880\\
10k/9k&5040&1000&--&0.436&1486.32&0.420&525.10&0.35x&1.251\\
10k/9k&5040&5000&128&0.428&1514.14&2.163&509.61&0.34x&2.159\\
10k/9k&10200&5000&128&0.985&1331.96&2.125&518.77&0.39x&3.105\\
10k/9k&50624&5000&128&4.256&1530.02&1.846&597.08&0.39x&10.874\\
10k/9k&100488&5000&128&8.595&1503.73&2.216&497.52&0.33x&14.339\\
30k/36k&120&5000&128&0.034&1821.49&5.039&868.47&0.48x&0.281\\
30k/36k&440&5000&128&0.125&1795.60&4.979&878.95&0.49x&1.187\\
30k/36k&960&5000&128&16.623&29.49&237.872&18.40&0.62x&2.461\\
30k/36k&10200&5000&128&39.130&133.09&56.585&77.34&0.58x&7.149\\
30k/36k&50624&5000&128&124.844&207.04&35.913&121.86&0.59x&23.372\\
30k/36k&100488&5000&128&268.417&191.15&38.993&112.24&0.59x&30.913\\
100k/94k&10200&5000&128&7.789&1730.33&10.031&1128.98&0.65x&27.531\\
100k/94k&50624&100&100&37.299&1793.31&0.211&1074.85&0.60x&1.400\\
100k/94k&50624&500&32&36.661&1824.52&3.022&374.70&0.21x&2.559\\
100k/94k&50624&500&128&37.288&1793.87&0.982&1153.24&0.64x&4.864\\
100k/94k&50624&500&256&39.112&1710.18&0.944&1199.83&0.70x&6.798\\
100k/94k&50624&1000&128&37.608&1778.59&1.865&1214.39&0.68x&8.335\\
100k/94k&50624&5000&--&35.814&1867.66&9.579&1182.32&0.63x&119.365\\
100k/94k&50624&5000&128&37.333&1791.67&9.528&1188.58&0.66x&36.243\\
100k/94k&50624&10000&128&36.716&1821.81&19.334&1171.55&0.64x&88.920\\
100k/94k&50624&50000&128&37.176&1799.24&96.469&1173.99&0.65x&350.826\\
100k/94k&50624&100000&128&36.404&1837.43&190.733&1187.56&0.65x&685.353\\
100k/94k&100488&5000&--&72.743&1825.25&10.039&1128.19&0.62x&230.564\\
100k/94k&100488&10000&128&120.953&1097.73&19.598&1155.78&1.05x&137.326\\
\bottomrule
\end{tabular}}

大网格下 SI 的 $E_{\rm SI}$ 为约 1.1--1.9 Gcell-sweeps/s，RSI 的 $E_{\rm RSI}$ 通常为约 1.1--1.2 Gcell-sweeps/s，因此在相同 sweep 工作单位下 RSI 的并行效率为 SI 的 0.58--0.70x；$B=32$ 会进一步降至 0.21x。$S_{316}$、10,000 样本点为 1.05x，属于测量波动下的接近持平。RSI 的端到端优势来自减少必须计算的方向/样本总工作量，而不是单个 sweep 比 SI 更高效。代表性 NCU profile：256 threads/block、achieved occupancy 29.84\%、L2 hit 96.47\%、L1/TEX hit 27.68\%、DRAM throughput 3.05\% peak。

\subsect{8.4 网格、角度、样本与 batch 因子的 270 组测试}

本节原计划覆盖三种空间网格（10k、30k、100k）、六档方向数（120、440、960、5,040、10,200、50,624）、五档样本数（8、64、512、1,024、4,096）和三档请求 batch（64、128、256），共 $3\times6\times5\times3=270$ 组。当前完成 269 组：10k 与 30k 各 90 组，100k 完成 89 组；唯一缺少 100k/$S_{224}$/4,096 样本/请求 $B=256$。每张表按方向数、样本数、请求 batch 由少到多排列；实际 B 为运行时实际容量，小于请求 B 时表示样本数本身更小。

每种空间网格均使用本网格上的确定性 SI 场作为参考解，均固定为 $S_{316}$（100,488 个方向）、8,192 样本、请求/实际 $B=128$；其中样本数和 batch 仅为完整 Figure 5 运行配置，参考场本身取 SI fine 输出。三个参考 case 如下：

{\scriptsize\setlength{\tabcolsep}{4pt}
\begin{tabular}{lrl}
\toprule
空间网格/cells & 参考 SI 配置 & case\\
\midrule
10k/9,187 & $S_{316}$ / 100,488 directions / 8,192 / $B=128$ & \code{10k\_s316\_r8192\_b128\_rep1}\\
30k/36,470 & $S_{316}$ / 100,488 directions / 8,192 / $B=128$ & \code{30k\_s316\_r8192\_b128\_rep1}\\
100k/94,378 & $S_{316}$ / 100,488 directions / 8,192 / $B=128$ & \code{100k\_s316\_r8192\_b128\_rep1}\\
\bottomrule
\end{tabular}}

因此，每个 RSI-tail error 都是在相同空间网格、相同 cell 顺序上相对该网格 $S_{316}$ SI 场计算的相对 $L_2$ 误差；不再进行跨网格插值或最近邻映射。error 可直接用于解释 RSI 的随机样本、角度离散和 batch 因子，而不混入空间网格差异。

两种方法的并行效率使用同一 directional cell-sweep 单位（Mcell-sweeps/s）：
\[
E_{\rm SI}=\frac{\mathrm{cells}\times N_{\Omega}\times N}{t_{\rm SI,sweep}\times10^6},\qquad
E_{\rm RSI}=\frac{\mathrm{cells}\times N_s\times(N+10)}{t_{\rm RSI,sweep}\times10^6}.
\]
其中 $N_{\Omega}$ 是方向数，$N_s$ 是样本数，$N$ 是 SI 收敛迭代数。效率表中的 RSI/SI 是 $E_{\rm RSI}/E_{\rm SI}$，即同一 sweep 工作单位下 RSI 相对 SI 的并行吞吐提升倍数；时间表中的 RSI/SI 则为 $t_{\rm SI,total}/t_{\rm RSI,total}$，即端到端加速比。

{\tiny\setlength{\tabcolsep}{2.3pt}
\begin{longtable}{rrrrrrrr}
\caption{10k 网格（9187 cells）：端到端时间与同网格 S316 参考误差}\label{tab:factor-time-10k}\\
\toprule
方向数 & 样本数 & 请求 B & 实际 B & SI total/s & RSI total/s & RSI/SI & error (\%)\\
\midrule
\endfirsthead
\multicolumn{8}{l}{\small 续表：10k 网格（9187 cells）：端到端时间与同网格 S316 参考误差}\\
\toprule
方向数 & 样本数 & 请求 B & 实际 B & SI total/s & RSI total/s & RSI/SI & error (\%)\\
\midrule
\endhead
\bottomrule
\endfoot
120 & 8 & 64 & 8 & 0.037 & 0.057 & 0.65x & 28.411\\
120 & 8 & 128 & 8 & 0.026 & 0.048 & 0.54x & 28.411\\
120 & 8 & 256 & 8 & 0.028 & 0.067 & 0.42x & 28.411\\
120 & 64 & 64 & 64 & 0.023 & 0.043 & 0.54x & 15.044\\
120 & 64 & 128 & 64 & 0.026 & 0.047 & 0.55x & 15.044\\
120 & 64 & 256 & 64 & 0.023 & 0.045 & 0.51x & 15.044\\
120 & 512 & 64 & 64 & 0.028 & 0.396 & 0.07x & 12.985\\
120 & 512 & 128 & 128 & 0.024 & 0.272 & 0.09x & 12.985\\
120 & 512 & 256 & 256 & 0.027 & 0.117 & 0.23x & 12.985\\
120 & 1024 & 64 & 64 & 0.034 & 0.893 & 0.04x & 12.784\\
120 & 1024 & 128 & 128 & 0.039 & 0.517 & 0.07x & 12.784\\
120 & 1024 & 256 & 256 & 0.026 & 0.285 & 0.09x & 12.784\\
120 & 4096 & 64 & 64 & 0.031 & 3.534 & 0.01x & 12.608\\
120 & 4096 & 128 & 128 & 0.027 & 1.643 & 0.02x & 12.608\\
120 & 4096 & 256 & 256 & 0.033 & 1.122 & 0.03x & 12.608\\
440 & 8 & 64 & 8 & 0.117 & 0.076 & 1.54x & 19.145\\
440 & 8 & 128 & 8 & 0.103 & 0.069 & 1.50x & 19.145\\
440 & 8 & 256 & 8 & 0.077 & 0.055 & 1.39x & 19.145\\
440 & 64 & 64 & 64 & 0.091 & 0.112 & 0.81x & 5.821\\
440 & 64 & 128 & 64 & 0.081 & 0.093 & 0.87x & 5.821\\
440 & 64 & 256 & 64 & 0.087 & 0.120 & 0.72x & 5.821\\
440 & 512 & 64 & 64 & 0.078 & 0.568 & 0.14x & 8.483\\
440 & 512 & 128 & 128 & 0.085 & 0.263 & 0.32x & 8.483\\
440 & 512 & 256 & 256 & 0.075 & 0.228 & 0.33x & 8.483\\
440 & 1024 & 64 & 64 & 0.091 & 1.043 & 0.09x & 7.593\\
440 & 1024 & 128 & 128 & 0.083 & 0.533 & 0.16x & 7.593\\
440 & 1024 & 256 & 256 & 0.075 & 0.312 & 0.24x & 7.593\\
440 & 4096 & 64 & 64 & 0.075 & 4.004 & 0.02x & 6.395\\
440 & 4096 & 128 & 128 & 0.080 & 2.255 & 0.04x & 6.395\\
440 & 4096 & 256 & 256 & 0.097 & 1.199 & 0.08x & 6.395\\
960 & 8 & 64 & 8 & 0.201 & 0.078 & 2.58x & 30.489\\
960 & 8 & 128 & 8 & 0.153 & 0.090 & 1.70x & 30.489\\
960 & 8 & 256 & 8 & 0.158 & 0.053 & 3.00x & 30.489\\
960 & 64 & 64 & 64 & 0.159 & 0.179 & 0.89x & 8.910\\
960 & 64 & 128 & 64 & 0.164 & 0.132 & 1.24x & 8.910\\
960 & 64 & 256 & 64 & 0.174 & 0.117 & 1.48x & 8.910\\
960 & 512 & 64 & 64 & 0.156 & 0.579 & 0.27x & 6.091\\
960 & 512 & 128 & 128 & 0.149 & 0.338 & 0.44x & 6.091\\
960 & 512 & 256 & 256 & 0.147 & 0.248 & 0.59x & 6.091\\
960 & 1024 & 64 & 64 & 0.150 & 1.039 & 0.14x & 5.397\\
960 & 1024 & 128 & 128 & 0.155 & 0.647 & 0.24x & 5.397\\
960 & 1024 & 256 & 256 & 0.165 & 0.394 & 0.42x & 5.397\\
960 & 4096 & 64 & 64 & 0.150 & 4.244 & 0.04x & 4.598\\
960 & 4096 & 128 & 128 & 0.146 & 2.468 & 0.06x & 4.598\\
960 & 4096 & 256 & 256 & 0.147 & 1.437 & 0.10x & 4.598\\
5040 & 8 & 64 & 8 & 1.010 & 0.083 & 12.21x & 21.737\\
5040 & 8 & 128 & 8 & 0.704 & 0.051 & 13.80x & 21.737\\
5040 & 8 & 256 & 8 & 0.739 & 0.086 & 8.55x & 21.737\\
5040 & 64 & 64 & 64 & 0.703 & 0.249 & 2.82x & 9.197\\
5040 & 64 & 128 & 64 & 0.715 & 0.211 & 3.38x & 9.197\\
5040 & 64 & 256 & 64 & 0.720 & 0.141 & 5.10x & 9.197\\
5040 & 512 & 64 & 64 & 0.703 & 0.937 & 0.75x & 3.581\\
5040 & 512 & 128 & 128 & 0.723 & 0.693 & 1.04x & 3.581\\
5040 & 512 & 256 & 256 & 0.737 & 0.667 & 1.10x & 3.581\\
5040 & 1024 & 64 & 64 & 0.718 & 1.458 & 0.49x & 3.912\\
5040 & 1024 & 128 & 128 & 0.739 & 1.068 & 0.69x & 3.912\\
5040 & 1024 & 256 & 256 & 0.753 & 0.952 & 0.79x & 3.912\\
5040 & 4096 & 64 & 64 & 0.768 & 5.203 & 0.15x & 2.023\\
5040 & 4096 & 128 & 128 & 0.718 & 3.272 & 0.22x & 2.023\\
5040 & 4096 & 256 & 256 & 0.691 & 2.523 & 0.27x & 2.023\\
10200 & 8 & 64 & 8 & 1.878 & 0.147 & 12.76x & 23.893\\
10200 & 8 & 128 & 8 & 1.392 & 0.053 & 26.07x & 23.893\\
10200 & 8 & 256 & 8 & 1.462 & 0.075 & 19.45x & 23.893\\
10200 & 64 & 64 & 64 & 1.449 & 0.289 & 5.01x & 7.635\\
10200 & 64 & 128 & 64 & 1.408 & 0.207 & 6.80x & 7.635\\
10200 & 64 & 256 & 64 & 1.440 & 0.205 & 7.03x & 7.635\\
10200 & 512 & 64 & 64 & 1.526 & 1.182 & 1.29x & 4.252\\
10200 & 512 & 128 & 128 & 1.483 & 0.768 & 1.93x & 4.252\\
10200 & 512 & 256 & 256 & 1.510 & 0.759 & 1.99x & 4.252\\
10200 & 1024 & 64 & 64 & 1.484 & 1.485 & 1.00x & 2.145\\
10200 & 1024 & 128 & 128 & 1.438 & 1.236 & 1.16x & 2.145\\
10200 & 1024 & 256 & 256 & 1.492 & 1.300 & 1.15x & 2.145\\
10200 & 4096 & 64 & 64 & 1.461 & 5.010 & 0.29x & 1.738\\
10200 & 4096 & 128 & 128 & 1.398 & 4.305 & 0.32x & 1.738\\
10200 & 4096 & 256 & 256 & 1.405 & 3.609 & 0.39x & 1.738\\
50624 & 8 & 64 & 8 & 9.121 & 0.096 & 94.76x & 40.820\\
50624 & 8 & 128 & 8 & 6.980 & 0.088 & 79.26x & 40.820\\
50624 & 8 & 256 & 8 & 6.637 & 0.075 & 88.19x & 40.820\\
50624 & 64 & 64 & 64 & 6.725 & 0.293 & 22.93x & 21.174\\
50624 & 64 & 128 & 64 & 6.898 & 0.277 & 24.87x & 21.174\\
50624 & 64 & 256 & 64 & 7.049 & 0.289 & 24.41x & 21.174\\
50624 & 512 & 64 & 64 & 6.864 & 1.453 & 4.72x & 5.468\\
50624 & 512 & 128 & 128 & 7.019 & 1.086 & 6.46x & 5.468\\
50624 & 512 & 256 & 256 & 6.987 & 1.129 & 6.19x & 5.468\\
50624 & 1024 & 64 & 64 & 6.646 & 2.210 & 3.01x & 2.349\\
50624 & 1024 & 128 & 128 & 6.649 & 1.804 & 3.69x & 2.349\\
50624 & 1024 & 256 & 256 & 6.691 & 2.135 & 3.13x & 2.349\\
50624 & 4096 & 64 & 64 & 6.853 & 6.977 & 0.98x & 1.172\\
50624 & 4096 & 128 & 128 & 6.759 & 5.822 & 1.16x & 1.172\\
50624 & 4096 & 256 & 256 & 6.990 & 5.908 & 1.18x & 1.172\\
\end{longtable}
}

{\tiny\setlength{\tabcolsep}{2.3pt}
\begin{longtable}{rrrrrrrr}
\caption{30k 网格（36470 cells）：端到端时间与同网格 S316 参考误差}\label{tab:factor-time-30k}\\
\toprule
方向数 & 样本数 & 请求 B & 实际 B & SI total/s & RSI total/s & RSI/SI & error (\%)\\
\midrule
\endfirsthead
\multicolumn{8}{l}{\small 续表：30k 网格（36470 cells）：端到端时间与同网格 S316 参考误差}\\
\toprule
方向数 & 样本数 & 请求 B & 实际 B & SI total/s & RSI total/s & RSI/SI & error (\%)\\
\midrule
\endhead
\bottomrule
\endfoot
120 & 8 & 64 & 8 & 0.105 & 0.137 & 0.77x & 23.443\\
120 & 8 & 128 & 8 & 0.077 & 0.126 & 0.61x & 23.443\\
120 & 8 & 256 & 8 & 0.078 & 0.101 & 0.77x & 23.443\\
120 & 64 & 64 & 64 & 0.084 & 0.107 & 0.78x & 12.508\\
120 & 64 & 128 & 64 & 0.080 & 0.096 & 0.83x & 12.508\\
120 & 64 & 256 & 64 & 0.083 & 0.114 & 0.73x & 12.508\\
120 & 512 & 64 & 64 & 0.079 & 0.933 & 0.09x & 11.663\\
120 & 512 & 128 & 128 & 0.079 & 0.582 & 0.14x & 11.663\\
120 & 512 & 256 & 256 & 0.079 & 0.403 & 0.20x & 11.663\\
120 & 1024 & 64 & 64 & 0.077 & 1.583 & 0.05x & 11.157\\
120 & 1024 & 128 & 128 & 0.080 & 1.184 & 0.07x & 11.157\\
120 & 1024 & 256 & 256 & 0.078 & 0.892 & 0.09x & 11.157\\
120 & 4096 & 64 & 64 & 0.088 & 7.206 & 0.01x & 11.239\\
120 & 4096 & 128 & 128 & 0.082 & 4.567 & 0.02x & 11.239\\
120 & 4096 & 256 & 256 & 0.088 & 3.244 & 0.03x & 11.239\\
440 & 8 & 64 & 8 & 0.343 & 0.171 & 2.00x & 19.501\\
440 & 8 & 128 & 8 & 0.250 & 0.130 & 1.92x & 19.501\\
440 & 8 & 256 & 8 & 0.235 & 0.204 & 1.15x & 19.501\\
440 & 64 & 64 & 64 & 0.236 & 0.260 & 0.91x & 5.167\\
440 & 64 & 128 & 64 & 0.234 & 0.205 & 1.14x & 5.167\\
440 & 64 & 256 & 64 & 0.232 & 0.229 & 1.01x & 5.167\\
440 & 512 & 64 & 64 & 0.237 & 1.219 & 0.19x & 8.865\\
440 & 512 & 128 & 128 & 0.246 & 0.655 & 0.38x & 8.865\\
440 & 512 & 256 & 256 & 0.268 & 0.560 & 0.48x & 8.865\\
440 & 1024 & 64 & 64 & 0.237 & 2.237 & 0.11x & 7.838\\
440 & 1024 & 128 & 128 & 0.236 & 1.382 & 0.17x & 7.838\\
440 & 1024 & 256 & 256 & 0.250 & 1.009 & 0.25x & 7.838\\
440 & 4096 & 64 & 64 & 0.248 & 8.447 & 0.03x & 5.941\\
440 & 4096 & 128 & 128 & 0.250 & 5.280 & 0.05x & 5.941\\
440 & 4096 & 256 & 256 & 0.237 & 3.682 & 0.06x & 5.941\\
960 & 8 & 64 & 8 & 16.880 & 1.307 & 12.91x & 29.438\\
960 & 8 & 128 & 8 & 16.281 & 1.227 & 13.27x & 29.438\\
960 & 8 & 256 & 8 & 16.621 & 1.216 & 13.67x & 29.438\\
960 & 64 & 64 & 64 & 16.401 & 3.566 & 4.60x & 8.515\\
960 & 64 & 128 & 64 & 16.038 & 3.555 & 4.51x & 8.515\\
960 & 64 & 256 & 64 & 16.303 & 3.537 & 4.61x & 8.515\\
960 & 512 & 64 & 64 & 16.497 & 31.384 & 0.53x & 5.972\\
960 & 512 & 128 & 128 & 17.575 & 28.706 & 0.61x & 5.972\\
960 & 512 & 256 & 256 & 16.596 & 24.033 & 0.69x & 5.972\\
960 & 1024 & 64 & 64 & 16.579 & 54.950 & 0.30x & 5.301\\
960 & 1024 & 128 & 128 & 16.507 & 48.553 & 0.34x & 5.301\\
960 & 1024 & 256 & 256 & 16.615 & 44.225 & 0.38x & 5.301\\
960 & 4096 & 64 & 64 & 16.639 & 226.747 & 0.07x & 4.116\\
960 & 4096 & 128 & 128 & 17.037 & 204.750 & 0.08x & 4.116\\
960 & 4096 & 256 & 256 & 16.786 & 180.022 & 0.09x & 4.116\\
5040 & 8 & 64 & 8 & 36.049 & 0.241 & 149.75x & 18.626\\
5040 & 8 & 128 & 8 & 35.121 & 0.190 & 185.22x & 18.626\\
5040 & 8 & 256 & 8 & 35.955 & 0.164 & 218.65x & 18.626\\
5040 & 64 & 64 & 64 & 35.877 & 3.137 & 11.44x & 7.850\\
5040 & 64 & 128 & 64 & 35.636 & 2.800 & 12.73x & 7.850\\
5040 & 64 & 256 & 64 & 35.889 & 3.004 & 11.95x & 7.850\\
5040 & 512 & 64 & 64 & 36.049 & 22.849 & 1.58x & 3.764\\
5040 & 512 & 128 & 128 & 35.908 & 21.371 & 1.68x & 3.764\\
5040 & 512 & 256 & 256 & 36.381 & 20.016 & 1.82x & 3.764\\
5040 & 1024 & 64 & 64 & 37.680 & 29.969 & 1.26x & 4.013\\
5040 & 1024 & 128 & 128 & 35.706 & 29.148 & 1.22x & 4.013\\
5040 & 1024 & 256 & 256 & 35.536 & 25.849 & 1.37x & 4.013\\
5040 & 4096 & 64 & 64 & 35.357 & 104.329 & 0.34x & 1.945\\
5040 & 4096 & 128 & 128 & 34.532 & 100.030 & 0.35x & 1.945\\
5040 & 4096 & 256 & 256 & 34.855 & 88.341 & 0.39x & 1.945\\
10200 & 8 & 64 & 8 & 38.344 & 0.248 & 154.68x & 21.487\\
10200 & 8 & 128 & 8 & 36.628 & 0.138 & 264.93x & 21.487\\
10200 & 8 & 256 & 8 & 36.523 & 0.177 & 206.83x & 21.487\\
10200 & 64 & 64 & 64 & 36.662 & 0.638 & 57.43x & 6.297\\
10200 & 64 & 128 & 64 & 37.094 & 0.584 & 63.49x & 6.297\\
10200 & 64 & 256 & 64 & 36.780 & 0.544 & 67.60x & 6.297\\
10200 & 512 & 64 & 64 & 36.668 & 13.214 & 2.77x & 3.290\\
10200 & 512 & 128 & 128 & 43.879 & 10.197 & 4.30x & 3.290\\
10200 & 512 & 256 & 256 & 43.092 & 10.961 & 3.93x & 3.290\\
10200 & 1024 & 64 & 64 & 39.247 & 16.816 & 2.33x & 1.639\\
10200 & 1024 & 128 & 128 & 36.959 & 14.127 & 2.62x & 1.639\\
10200 & 1024 & 256 & 256 & 36.320 & 13.688 & 2.65x & 1.639\\
10200 & 4096 & 64 & 64 & 37.022 & 54.343 & 0.68x & 1.602\\
10200 & 4096 & 128 & 128 & 36.028 & 49.800 & 0.72x & 1.602\\
10200 & 4096 & 256 & 256 & 36.001 & 47.395 & 0.76x & 1.602\\
50624 & 8 & 64 & 8 & 140.023 & 0.212 & 659.60x & 35.121\\
50624 & 8 & 128 & 8 & 135.960 & 0.179 & 759.44x & 35.121\\
50624 & 8 & 256 & 8 & 130.253 & 0.108 & 1201.17x & 35.121\\
50624 & 64 & 64 & 64 & 129.774 & 2.792 & 46.49x & 19.619\\
50624 & 64 & 128 & 64 & 129.525 & 2.687 & 48.21x & 19.619\\
50624 & 64 & 256 & 64 & 129.269 & 2.750 & 47.00x & 19.619\\
50624 & 512 & 64 & 64 & 161.923 & 12.237 & 13.23x & 4.767\\
50624 & 512 & 128 & 128 & 166.438 & 10.768 & 15.46x & 4.767\\
50624 & 512 & 256 & 256 & 169.810 & 9.654 & 17.59x & 4.767\\
50624 & 1024 & 64 & 64 & 164.983 & 14.386 & 11.47x & 2.395\\
50624 & 1024 & 128 & 128 & 155.774 & 15.516 & 10.04x & 2.395\\
50624 & 1024 & 256 & 256 & 136.740 & 13.874 & 9.86x & 2.395\\
50624 & 4096 & 64 & 64 & 154.865 & 48.708 & 3.18x & 1.054\\
50624 & 4096 & 128 & 128 & 168.836 & 45.147 & 3.74x & 1.054\\
50624 & 4096 & 256 & 256 & 145.068 & 46.633 & 3.11x & 1.054\\
\end{longtable}
}

{\tiny\setlength{\tabcolsep}{2.3pt}
\begin{longtable}{rrrrrrrr}
\caption{100k 网格（94378 cells）：端到端时间与同网格 S316 参考误差；缺少 S224/4096/请求 B=256}\label{tab:factor-time-100k}\\
\toprule
方向数 & 样本数 & 请求 B & 实际 B & SI total/s & RSI total/s & RSI/SI & error (\%)\\
\midrule
\endfirsthead
\multicolumn{8}{l}{\small 续表：100k 网格（94378 cells）：端到端时间与同网格 S316 参考误差；缺少 S224/4096/请求 B=256}\\
\toprule
方向数 & 样本数 & 请求 B & 实际 B & SI total/s & RSI total/s & RSI/SI & error (\%)\\
\midrule
\endhead
\bottomrule
\endfoot
120 & 8 & 64 & 8 & 0.539 & 0.459 & 1.17x & 30.574\\
120 & 8 & 128 & 8 & 0.205 & 0.217 & 0.95x & 30.574\\
120 & 8 & 256 & 8 & 0.196 & 0.205 & 0.96x & 30.574\\
120 & 64 & 64 & 64 & 0.188 & 0.227 & 0.83x & 16.386\\
120 & 64 & 128 & 64 & 0.189 & 0.213 & 0.89x & 16.386\\
120 & 64 & 256 & 64 & 0.187 & 0.214 & 0.87x & 16.386\\
120 & 512 & 64 & 64 & 0.193 & 1.486 & 0.13x & 14.481\\
120 & 512 & 128 & 128 & 0.184 & 1.254 & 0.15x & 14.481\\
120 & 512 & 256 & 256 & 0.186 & 0.950 & 0.20x & 14.481\\
120 & 1024 & 64 & 64 & 0.181 & 3.251 & 0.06x & 14.130\\
120 & 1024 & 128 & 128 & 0.184 & 2.425 & 0.08x & 14.130\\
120 & 1024 & 256 & 256 & 0.182 & 1.906 & 0.10x & 14.130\\
120 & 4096 & 64 & 64 & 0.184 & 13.437 & 0.01x & 13.984\\
120 & 4096 & 128 & 128 & 0.185 & 9.472 & 0.02x & 13.984\\
120 & 4096 & 256 & 256 & 0.195 & 7.662 & 0.03x & 13.984\\
440 & 8 & 64 & 8 & 1.446 & 0.553 & 2.62x & 23.194\\
440 & 8 & 128 & 8 & 0.622 & 0.248 & 2.51x & 23.194\\
440 & 8 & 256 & 8 & 0.610 & 0.365 & 1.67x & 23.194\\
440 & 64 & 64 & 64 & 0.639 & 1.267 & 0.50x & 6.777\\
440 & 64 & 128 & 64 & 0.629 & 0.423 & 1.49x & 6.777\\
440 & 64 & 256 & 64 & 0.632 & 0.412 & 1.53x & 6.777\\
440 & 512 & 64 & 64 & 0.591 & 2.060 & 0.29x & 9.297\\
440 & 512 & 128 & 128 & 0.594 & 1.527 & 0.39x & 9.297\\
440 & 512 & 256 & 256 & 0.625 & 1.230 & 0.51x & 9.297\\
440 & 1024 & 64 & 64 & 0.597 & 4.563 & 0.13x & 8.364\\
440 & 1024 & 128 & 128 & 0.693 & 3.246 & 0.21x & 8.364\\
440 & 1024 & 256 & 256 & 0.689 & 2.683 & 0.26x & 8.364\\
440 & 4096 & 64 & 64 & 0.672 & 17.722 & 0.04x & 7.006\\
440 & 4096 & 128 & 128 & 0.626 & 10.877 & 0.06x & 7.006\\
440 & 4096 & 256 & 256 & 0.598 & 8.772 & 0.07x & 7.006\\
960 & 8 & 64 & 8 & 3.789 & 0.723 & 5.24x & 33.702\\
960 & 8 & 128 & 8 & 1.438 & 0.298 & 4.82x & 33.702\\
960 & 8 & 256 & 8 & 1.484 & 0.313 & 4.75x & 33.702\\
960 & 64 & 64 & 64 & 1.434 & 2.861 & 0.50x & 10.201\\
960 & 64 & 128 & 64 & 1.463 & 0.798 & 1.83x & 10.201\\
960 & 64 & 256 & 64 & 1.665 & 0.807 & 2.06x & 10.201\\
960 & 512 & 64 & 64 & 1.560 & 3.693 & 0.42x & 6.815\\
960 & 512 & 128 & 128 & 1.447 & 2.362 & 0.61x & 6.815\\
960 & 512 & 256 & 256 & 1.433 & 1.916 & 0.75x & 6.815\\
960 & 1024 & 64 & 64 & 1.419 & 5.662 & 0.25x & 5.989\\
960 & 1024 & 128 & 128 & 1.432 & 4.240 & 0.34x & 5.989\\
960 & 1024 & 256 & 256 & 1.383 & 3.250 & 0.43x & 5.989\\
960 & 4096 & 64 & 64 & 1.406 & 21.263 & 0.07x & 5.025\\
960 & 4096 & 128 & 128 & 1.271 & 14.328 & 0.09x & 5.025\\
960 & 4096 & 256 & 256 & 1.247 & 10.128 & 0.12x & 5.025\\
5040 & 8 & 64 & 8 & 17.449 & 0.662 & 26.35x & 24.457\\
5040 & 8 & 128 & 8 & 6.325 & 0.217 & 29.15x & 24.457\\
5040 & 8 & 256 & 8 & 5.955 & 0.293 & 20.33x & 24.457\\
5040 & 64 & 64 & 64 & 5.873 & 3.966 & 1.48x & 9.633\\
5040 & 64 & 128 & 64 & 5.982 & 1.465 & 4.08x & 9.633\\
5040 & 64 & 256 & 64 & 6.056 & 1.556 & 3.89x & 9.633\\
5040 & 512 & 64 & 64 & 6.201 & 12.680 & 0.49x & 4.001\\
5040 & 512 & 128 & 128 & 6.116 & 4.837 & 1.26x & 4.001\\
5040 & 512 & 256 & 256 & 6.303 & 4.960 & 1.27x & 4.001\\
5040 & 1024 & 64 & 64 & 6.276 & 11.034 & 0.57x & 4.237\\
5040 & 1024 & 128 & 128 & 6.093 & 7.925 & 0.77x & 4.237\\
5040 & 1024 & 256 & 256 & 6.889 & 7.628 & 0.90x & 4.237\\
5040 & 4096 & 64 & 64 & 6.090 & 30.928 & 0.20x & 2.199\\
5040 & 4096 & 128 & 128 & 6.005 & 30.395 & 0.20x & 2.199\\
5040 & 4096 & 256 & 256 & 7.348 & 25.028 & 0.29x & 2.199\\
10200 & 8 & 64 & 8 & 14.849 & 0.775 & 19.15x & 26.621\\
10200 & 8 & 128 & 8 & 13.038 & 0.256 & 50.97x & 26.621\\
10200 & 8 & 256 & 8 & 13.012 & 0.283 & 45.98x & 26.621\\
10200 & 64 & 64 & 64 & 12.515 & 1.188 & 10.54x & 8.436\\
10200 & 64 & 128 & 64 & 12.032 & 1.290 & 9.33x & 8.436\\
10200 & 64 & 256 & 64 & 11.438 & 1.257 & 9.10x & 8.436\\
10200 & 512 & 64 & 64 & 11.428 & 5.464 & 2.09x & 4.473\\
10200 & 512 & 128 & 128 & 11.639 & 5.070 & 2.30x & 4.473\\
10200 & 512 & 256 & 256 & 12.729 & 7.408 & 1.72x & 4.473\\
10200 & 1024 & 64 & 64 & 14.962 & 11.548 & 1.30x & 2.231\\
10200 & 1024 & 128 & 128 & 13.264 & 11.419 & 1.16x & 2.231\\
10200 & 1024 & 256 & 256 & 13.074 & 9.155 & 1.43x & 2.231\\
10200 & 4096 & 64 & 64 & 13.316 & 34.837 & 0.38x & 1.841\\
10200 & 4096 & 128 & 128 & 12.460 & 34.334 & 0.36x & 1.841\\
10200 & 4096 & 256 & 256 & 13.590 & 30.248 & 0.45x & 1.841\\
50624 & 8 & 64 & 8 & 134.551 & 0.812 & 165.61x & 43.011\\
50624 & 8 & 128 & 8 & 134.012 & 0.263 & 510.40x & 43.011\\
50624 & 8 & 256 & 8 & 127.902 & 0.252 & 507.72x & 43.011\\
50624 & 64 & 64 & 64 & 168.503 & 6.394 & 26.35x & 21.374\\
50624 & 64 & 128 & 64 & 144.489 & 1.439 & 100.44x & 21.374\\
50624 & 64 & 256 & 64 & 136.794 & 1.066 & 128.31x & 21.374\\
50624 & 512 & 64 & 64 & 136.577 & 6.017 & 22.70x & 5.573\\
50624 & 512 & 128 & 128 & 139.865 & 6.334 & 22.08x & 5.573\\
50624 & 512 & 256 & 256 & 141.051 & 8.060 & 17.50x & 5.573\\
50624 & 1024 & 64 & 64 & 133.688 & 10.468 & 12.77x & 2.562\\
50624 & 1024 & 128 & 128 & 139.067 & 12.018 & 11.57x & 2.562\\
50624 & 1024 & 256 & 256 & 143.106 & 13.725 & 10.43x & 2.562\\
50624 & 4096 & 64 & 64 & 138.122 & 44.853 & 3.08x & 1.252\\
50624 & 4096 & 128 & 128 & 142.833 & 45.893 & 3.11x & 1.252\\
\end{longtable}
}

{\tiny\setlength{\tabcolsep}{2.3pt}
\begin{longtable}{rrrrrrrrrr}
\caption{10k 网格（9187 cells）：SI/RSI 并行效率}\label{tab:factor-efficiency-10k}\\
\toprule
方向数 & 样本数 & 请求 B & 实际 B & SI sweep/s & $E_{\rm SI}$ & RSI sweep/s & $E_{\rm RSI}$ & RSI/SI & RSI plan/s\\
\midrule
\endfirsthead
\multicolumn{10}{l}{\small 续表：10k 网格（9187 cells）：SI/RSI 并行效率}\\
\toprule
方向数 & 样本数 & 请求 B & 实际 B & SI sweep/s & $E_{\rm SI}$ & RSI sweep/s & $E_{\rm RSI}$ & RSI/SI & RSI plan/s\\
\midrule
\endhead
\bottomrule
\endfoot
120 & 8 & 64 & 8 & 0.010 & 1498.6 & 0.035 & 50.1 & 0.03x & 0.018\\
120 & 8 & 128 & 8 & 0.011 & 1459.5 & 0.037 & 47.3 & 0.03x & 0.007\\
120 & 8 & 256 & 8 & 0.011 & 1414.9 & 0.055 & 32.2 & 0.02x & 0.007\\
120 & 64 & 64 & 64 & 0.010 & 1488.9 & 0.035 & 403.8 & 0.27x & 0.003\\
120 & 64 & 128 & 64 & 0.011 & 1436.2 & 0.038 & 373.7 & 0.26x & 0.004\\
120 & 64 & 256 & 64 & 0.010 & 1498.6 & 0.037 & 385.1 & 0.26x & 0.004\\
120 & 512 & 64 & 64 & 0.011 & 1375.0 & 0.342 & 330.1 & 0.24x & 0.019\\
120 & 512 & 128 & 128 & 0.010 & 1513.4 & 0.237 & 476.1 & 0.31x & 0.012\\
120 & 512 & 256 & 256 & 0.011 & 1459.7 & 0.097 & 1159.0 & 0.79x & 0.006\\
120 & 1024 & 64 & 64 & 0.013 & 1215.2 & 0.775 & 291.2 & 0.24x & 0.041\\
120 & 1024 & 128 & 128 & 0.011 & 1456.4 & 0.450 & 501.8 & 0.34x & 0.022\\
120 & 1024 & 256 & 256 & 0.011 & 1454.8 & 0.246 & 917.2 & 0.63x & 0.010\\
120 & 4096 & 64 & 64 & 0.011 & 1467.6 & 3.092 & 292.1 & 0.20x & 0.154\\
120 & 4096 & 128 & 128 & 0.011 & 1438.2 & 1.429 & 632.1 & 0.44x & 0.071\\
120 & 4096 & 256 & 256 & 0.010 & 1499.4 & 0.980 & 921.5 & 0.61x & 0.039\\
440 & 8 & 64 & 8 & 0.047 & 1200.6 & 0.042 & 42.3 & 0.04x & 0.029\\
440 & 8 & 128 & 8 & 0.046 & 1223.9 & 0.047 & 37.8 & 0.03x & 0.017\\
440 & 8 & 256 & 8 & 0.042 & 1340.2 & 0.039 & 44.9 & 0.03x & 0.011\\
440 & 64 & 64 & 64 & 0.048 & 1177.8 & 0.043 & 326.3 & 0.28x & 0.063\\
440 & 64 & 128 & 64 & 0.045 & 1256.4 & 0.049 & 290.3 & 0.23x & 0.037\\
440 & 64 & 256 & 64 & 0.048 & 1182.8 & 0.056 & 251.8 & 0.21x & 0.055\\
440 & 512 & 64 & 64 & 0.042 & 1347.1 & 0.447 & 252.3 & 0.19x & 0.076\\
440 & 512 & 128 & 128 & 0.047 & 1215.5 & 0.186 & 605.9 & 0.50x & 0.054\\
440 & 512 & 256 & 256 & 0.043 & 1321.1 & 0.170 & 665.1 & 0.50x & 0.039\\
440 & 1024 & 64 & 64 & 0.049 & 1157.1 & 0.837 & 269.7 & 0.23x & 0.123\\
440 & 1024 & 128 & 128 & 0.047 & 1204.9 & 0.413 & 546.8 & 0.45x & 0.077\\
440 & 1024 & 256 & 256 & 0.042 & 1339.9 & 0.221 & 1022.1 & 0.76x & 0.063\\
440 & 4096 & 64 & 64 & 0.042 & 1353.1 & 3.278 & 275.5 & 0.20x & 0.406\\
440 & 4096 & 128 & 128 & 0.043 & 1322.2 & 1.837 & 491.7 & 0.37x & 0.232\\
440 & 4096 & 256 & 256 & 0.050 & 1137.5 & 0.963 & 938.1 & 0.82x & 0.124\\
960 & 8 & 64 & 8 & 0.087 & 1420.4 & 0.047 & 37.3 & 0.03x & 0.025\\
960 & 8 & 128 & 8 & 0.087 & 1412.5 & 0.067 & 26.5 & 0.02x & 0.017\\
960 & 8 & 256 & 8 & 0.091 & 1361.7 & 0.037 & 48.0 & 0.04x & 0.011\\
960 & 64 & 64 & 64 & 0.089 & 1391.5 & 0.056 & 250.9 & 0.18x & 0.115\\
960 & 64 & 128 & 64 & 0.095 & 1296.6 & 0.062 & 226.8 & 0.17x & 0.061\\
960 & 64 & 256 & 64 & 0.097 & 1277.1 & 0.051 & 277.9 & 0.22x & 0.059\\
960 & 512 & 64 & 64 & 0.091 & 1352.2 & 0.393 & 287.5 & 0.21x & 0.138\\
960 & 512 & 128 & 128 & 0.085 & 1449.1 & 0.201 & 561.3 & 0.39x & 0.110\\
960 & 512 & 256 & 256 & 0.086 & 1430.5 & 0.130 & 870.5 & 0.61x & 0.096\\
960 & 1024 & 64 & 64 & 0.086 & 1442.1 & 0.752 & 300.3 & 0.21x & 0.204\\
960 & 1024 & 128 & 128 & 0.091 & 1357.6 & 0.410 & 551.2 & 0.41x & 0.179\\
960 & 1024 & 256 & 256 & 0.093 & 1332.1 & 0.227 & 995.1 & 0.75x & 0.132\\
960 & 4096 & 64 & 64 & 0.087 & 1417.0 & 3.206 & 281.7 & 0.20x & 0.695\\
960 & 4096 & 128 & 128 & 0.085 & 1444.6 & 1.749 & 516.4 & 0.36x & 0.477\\
960 & 4096 & 256 & 256 & 0.084 & 1475.8 & 0.994 & 908.6 & 0.62x & 0.299\\
5040 & 8 & 64 & 8 & 0.418 & 1552.0 & 0.043 & 41.4 & 0.03x & 0.035\\
5040 & 8 & 128 & 8 & 0.446 & 1454.1 & 0.035 & 50.6 & 0.03x & 0.012\\
5040 & 8 & 256 & 8 & 0.457 & 1417.7 & 0.055 & 31.9 & 0.02x & 0.025\\
5040 & 64 & 64 & 64 & 0.427 & 1519.4 & 0.057 & 249.1 & 0.16x & 0.183\\
5040 & 64 & 128 & 64 & 0.454 & 1428.4 & 0.038 & 370.5 & 0.26x & 0.165\\
5040 & 64 & 256 & 64 & 0.450 & 1440.4 & 0.035 & 398.7 & 0.28x & 0.098\\
5040 & 512 & 64 & 64 & 0.432 & 1502.0 & 0.471 & 239.7 & 0.16x & 0.395\\
5040 & 512 & 128 & 128 & 0.445 & 1455.5 & 0.189 & 598.4 & 0.41x & 0.461\\
5040 & 512 & 256 & 256 & 0.437 & 1482.5 & 0.131 & 859.2 & 0.58x & 0.503\\
5040 & 1024 & 64 & 64 & 0.435 & 1489.5 & 0.936 & 241.2 & 0.16x & 0.387\\
5040 & 1024 & 128 & 128 & 0.455 & 1423.2 & 0.423 & 533.5 & 0.37x & 0.564\\
5040 & 1024 & 256 & 256 & 0.453 & 1431.2 & 0.294 & 767.6 & 0.54x & 0.596\\
5040 & 4096 & 64 & 64 & 0.452 & 1435.0 & 3.406 & 265.1 & 0.18x & 1.279\\
5040 & 4096 & 128 & 128 & 0.434 & 1494.4 & 1.659 & 544.4 & 0.36x & 1.278\\
5040 & 4096 & 256 & 256 & 0.426 & 1522.7 & 1.082 & 835.0 & 0.55x & 1.198\\
10200 & 8 & 64 & 8 & 0.867 & 1512.4 & 0.098 & 18.0 & 0.01x & 0.041\\
10200 & 8 & 128 & 8 & 0.864 & 1518.8 & 0.035 & 50.8 & 0.03x & 0.014\\
10200 & 8 & 256 & 8 & 0.867 & 1512.4 & 0.047 & 37.9 & 0.03x & 0.022\\
10200 & 64 & 64 & 64 & 0.882 & 1487.2 & 0.046 & 306.8 & 0.21x & 0.232\\
10200 & 64 & 128 & 64 & 0.872 & 1504.2 & 0.040 & 353.0 & 0.23x & 0.158\\
10200 & 64 & 256 & 64 & 0.876 & 1498.4 & 0.042 & 332.2 & 0.22x & 0.152\\
10200 & 512 & 64 & 64 & 0.898 & 1461.5 & 0.581 & 194.3 & 0.13x & 0.505\\
10200 & 512 & 128 & 128 & 0.883 & 1485.0 & 0.181 & 625.1 & 0.42x & 0.533\\
10200 & 512 & 256 & 256 & 0.909 & 1442.6 & 0.107 & 1056.8 & 0.73x & 0.608\\
10200 & 1024 & 64 & 64 & 0.887 & 1479.1 & 0.884 & 255.3 & 0.17x & 0.445\\
10200 & 1024 & 128 & 128 & 0.882 & 1486.7 & 0.458 & 493.2 & 0.33x & 0.669\\
10200 & 1024 & 256 & 256 & 0.916 & 1432.2 & 0.292 & 773.3 & 0.54x & 0.917\\
10200 & 4096 & 64 & 64 & 0.882 & 1487.3 & 3.034 & 297.7 & 0.20x & 1.410\\
10200 & 4096 & 128 & 128 & 0.862 & 1522.0 & 2.045 & 441.6 & 0.29x & 1.831\\
10200 & 4096 & 256 & 256 & 0.871 & 1505.8 & 1.049 & 861.3 & 0.57x & 2.228\\
50624 & 8 & 64 & 8 & 4.312 & 1510.1 & 0.052 & 33.9 & 0.02x & 0.037\\
50624 & 8 & 128 & 8 & 4.347 & 1497.8 & 0.059 & 30.0 & 0.02x & 0.022\\
50624 & 8 & 256 & 8 & 4.199 & 1550.5 & 0.051 & 34.3 & 0.02x & 0.017\\
50624 & 64 & 64 & 64 & 4.364 & 1492.0 & 0.036 & 388.0 & 0.26x & 0.237\\
50624 & 64 & 128 & 64 & 4.421 & 1472.9 & 0.037 & 378.0 & 0.26x & 0.221\\
50624 & 64 & 256 & 64 & 4.249 & 1532.3 & 0.047 & 299.6 & 0.20x & 0.223\\
50624 & 512 & 64 & 64 & 4.266 & 1526.3 & 0.443 & 255.1 & 0.17x & 0.866\\
50624 & 512 & 128 & 128 & 4.246 & 1533.5 & 0.179 & 629.7 & 0.41x & 0.786\\
50624 & 512 & 256 & 256 & 4.303 & 1513.3 & 0.117 & 967.7 & 0.64x & 0.892\\
50624 & 1024 & 64 & 64 & 4.229 & 1539.5 & 0.740 & 305.0 & 0.20x & 1.163\\
50624 & 1024 & 128 & 128 & 4.210 & 1546.6 & 0.408 & 553.0 & 0.36x & 1.145\\
50624 & 1024 & 256 & 256 & 4.241 & 1535.3 & 0.300 & 753.6 & 0.49x & 1.582\\
50624 & 4096 & 64 & 64 & 4.240 & 1535.7 & 3.218 & 280.6 & 0.18x & 2.555\\
50624 & 4096 & 128 & 128 & 4.335 & 1502.1 & 1.772 & 509.8 & 0.34x & 3.019\\
50624 & 4096 & 256 & 256 & 4.333 & 1502.6 & 1.091 & 827.7 & 0.55x & 3.831\\
\end{longtable}
}

{\tiny\setlength{\tabcolsep}{2.3pt}
\begin{longtable}{rrrrrrrrrr}
\caption{30k 网格（36470 cells）：SI/RSI 并行效率}\label{tab:factor-efficiency-30k}\\
\toprule
方向数 & 样本数 & 请求 B & 实际 B & SI sweep/s & $E_{\rm SI}$ & RSI sweep/s & $E_{\rm RSI}$ & RSI/SI & RSI plan/s\\
\midrule
\endfirsthead
\multicolumn{10}{l}{\small 续表：30k 网格（36470 cells）：SI/RSI 并行效率}\\
\toprule
方向数 & 样本数 & 请求 B & 实际 B & SI sweep/s & $E_{\rm SI}$ & RSI sweep/s & $E_{\rm RSI}$ & RSI/SI & RSI plan/s\\
\midrule
\endhead
\bottomrule
\endfoot
120 & 8 & 64 & 8 & 0.034 & 1798.5 & 0.072 & 97.8 & 0.05x & 0.059\\
120 & 8 & 128 & 8 & 0.033 & 1877.3 & 0.095 & 73.5 & 0.04x & 0.024\\
120 & 8 & 256 & 8 & 0.033 & 1878.5 & 0.068 & 103.3 & 0.05x & 0.027\\
120 & 64 & 64 & 64 & 0.033 & 1835.3 & 0.084 & 665.7 & 0.36x & 0.016\\
120 & 64 & 128 & 64 & 0.033 & 1834.0 & 0.075 & 748.7 & 0.41x & 0.015\\
120 & 64 & 256 & 64 & 0.034 & 1813.6 & 0.093 & 600.7 & 0.33x & 0.014\\
120 & 512 & 64 & 64 & 0.034 & 1776.4 & 0.822 & 545.4 & 0.31x & 0.059\\
120 & 512 & 128 & 128 & 0.034 & 1820.2 & 0.509 & 880.6 & 0.48x & 0.036\\
120 & 512 & 256 & 256 & 0.033 & 1852.4 & 0.350 & 1281.7 & 0.69x & 0.021\\
120 & 1024 & 64 & 64 & 0.033 & 1884.8 & 1.376 & 651.4 & 0.35x & 0.108\\
120 & 1024 & 128 & 128 & 0.033 & 1862.6 & 1.038 & 863.7 & 0.46x & 0.060\\
120 & 1024 & 256 & 256 & 0.033 & 1863.4 & 0.794 & 1128.7 & 0.61x & 0.033\\
120 & 4096 & 64 & 64 & 0.037 & 1634.6 & 6.341 & 565.4 & 0.35x & 0.415\\
120 & 4096 & 128 & 128 & 0.035 & 1753.1 & 4.013 & 893.3 & 0.51x & 0.219\\
120 & 4096 & 256 & 256 & 0.039 & 1567.5 & 2.870 & 1249.2 & 0.80x & 0.112\\
440 & 8 & 64 & 8 & 0.128 & 1750.9 & 0.072 & 97.6 & 0.06x & 0.091\\
440 & 8 & 128 & 8 & 0.130 & 1726.9 & 0.085 & 82.0 & 0.05x & 0.036\\
440 & 8 & 256 & 8 & 0.123 & 1824.3 & 0.156 & 44.8 & 0.02x & 0.037\\
440 & 64 & 64 & 64 & 0.122 & 1848.6 & 0.128 & 439.0 & 0.24x & 0.120\\
440 & 64 & 128 & 64 & 0.122 & 1834.8 & 0.090 & 624.1 & 0.34x & 0.105\\
440 & 64 & 256 & 64 & 0.123 & 1832.2 & 0.112 & 500.2 & 0.27x & 0.105\\
440 & 512 & 64 & 64 & 0.125 & 1792.4 & 0.885 & 506.6 & 0.28x & 0.249\\
440 & 512 & 128 & 128 & 0.128 & 1761.9 & 0.433 & 1034.4 & 0.59x & 0.174\\
440 & 512 & 256 & 256 & 0.143 & 1575.2 & 0.388 & 1154.6 & 0.73x & 0.132\\
440 & 1024 & 64 & 64 & 0.124 & 1814.1 & 1.669 & 537.1 & 0.30x & 0.410\\
440 & 1024 & 128 & 128 & 0.124 & 1808.1 & 1.025 & 874.4 & 0.48x & 0.254\\
440 & 1024 & 256 & 256 & 0.126 & 1780.9 & 0.755 & 1187.7 & 0.67x & 0.177\\
440 & 4096 & 64 & 64 & 0.131 & 1709.8 & 6.357 & 564.0 & 0.33x & 1.417\\
440 & 4096 & 128 & 128 & 0.130 & 1732.0 & 4.072 & 880.4 & 0.51x & 0.784\\
440 & 4096 & 256 & 256 & 0.125 & 1801.6 & 3.025 & 1185.3 & 0.66x & 0.417\\
960 & 8 & 64 & 8 & 15.928 & 30.8 & 1.168 & 6.0 & 0.19x & 0.104\\
960 & 8 & 128 & 8 & 15.800 & 31.0 & 1.151 & 6.1 & 0.20x & 0.042\\
960 & 8 & 256 & 8 & 16.013 & 30.6 & 1.140 & 6.1 & 0.20x & 0.041\\
960 & 64 & 64 & 64 & 15.801 & 31.0 & 3.275 & 17.1 & 0.55x & 0.198\\
960 & 64 & 128 & 64 & 15.448 & 31.7 & 3.286 & 17.0 & 0.54x & 0.176\\
960 & 64 & 256 & 64 & 15.715 & 31.2 & 3.267 & 17.1 & 0.55x & 0.176\\
960 & 512 & 64 & 64 & 15.901 & 30.8 & 30.127 & 14.9 & 0.48x & 0.425\\
960 & 512 & 128 & 128 & 16.937 & 28.9 & 27.614 & 16.2 & 0.56x & 0.346\\
960 & 512 & 256 & 256 & 15.982 & 30.7 & 23.151 & 19.4 & 0.63x & 0.262\\
960 & 1024 & 64 & 64 & 15.964 & 30.7 & 52.771 & 17.0 & 0.55x & 0.707\\
960 & 1024 & 128 & 128 & 15.901 & 30.8 & 46.763 & 19.2 & 0.62x & 0.511\\
960 & 1024 & 256 & 256 & 15.988 & 30.7 & 42.703 & 21.0 & 0.68x & 0.350\\
960 & 4096 & 64 & 64 & 16.004 & 30.6 & 218.280 & 16.4 & 0.54x & 2.393\\
960 & 4096 & 128 & 128 & 16.410 & 29.9 & 197.789 & 18.1 & 0.61x & 1.523\\
960 & 4096 & 256 & 256 & 16.148 & 30.4 & 174.429 & 20.6 & 0.68x & 0.881\\
5040 & 8 & 64 & 8 & 33.160 & 77.6 & 0.067 & 104.5 & 1.35x & 0.167\\
5040 & 8 & 128 & 8 & 33.000 & 78.0 & 0.074 & 95.0 & 1.22x & 0.108\\
5040 & 8 & 256 & 8 & 33.846 & 76.0 & 0.074 & 94.4 & 1.24x & 0.082\\
5040 & 64 & 64 & 64 & 33.892 & 75.9 & 2.274 & 24.6 & 0.32x & 0.792\\
5040 & 64 & 128 & 64 & 33.713 & 76.3 & 2.263 & 24.8 & 0.32x & 0.462\\
5040 & 64 & 256 & 64 & 33.936 & 75.8 & 2.384 & 23.5 & 0.31x & 0.537\\
5040 & 512 & 64 & 64 & 33.954 & 75.8 & 21.242 & 21.1 & 0.28x & 1.042\\
5040 & 512 & 128 & 128 & 34.026 & 75.6 & 19.522 & 23.0 & 0.30x & 1.207\\
5040 & 512 & 256 & 256 & 33.636 & 76.5 & 17.785 & 25.2 & 0.33x & 1.682\\
5040 & 1024 & 64 & 64 & 34.898 & 73.7 & 27.565 & 32.5 & 0.44x & 1.478\\
5040 & 1024 & 128 & 128 & 33.400 & 77.0 & 26.677 & 33.6 & 0.44x & 1.598\\
5040 & 1024 & 256 & 256 & 33.641 & 76.5 & 23.540 & 38.1 & 0.50x & 1.559\\
5040 & 4096 & 64 & 64 & 33.072 & 77.8 & 96.050 & 37.3 & 0.48x & 4.907\\
5040 & 4096 & 128 & 128 & 32.522 & 79.1 & 92.597 & 38.7 & 0.49x & 4.284\\
5040 & 4096 & 256 & 256 & 33.038 & 77.9 & 82.145 & 43.6 & 0.56x & 3.449\\
10200 & 8 & 64 & 8 & 34.128 & 152.6 & 0.107 & 65.4 & 0.43x & 0.132\\
10200 & 8 & 128 & 8 & 33.941 & 153.4 & 0.084 & 83.8 & 0.55x & 0.043\\
10200 & 8 & 256 & 8 & 33.995 & 153.2 & 0.118 & 59.3 & 0.39x & 0.050\\
10200 & 64 & 64 & 64 & 34.120 & 152.6 & 0.169 & 330.9 & 2.17x & 0.449\\
10200 & 64 & 128 & 64 & 34.509 & 150.9 & 0.144 & 388.9 & 2.58x & 0.420\\
10200 & 64 & 256 & 64 & 34.152 & 152.5 & 0.107 & 521.7 & 3.42x & 0.418\\
10200 & 512 & 64 & 64 & 34.076 & 152.8 & 8.892 & 50.4 & 0.33x & 3.938\\
10200 & 512 & 128 & 128 & 39.090 & 133.2 & 7.551 & 59.3 & 0.45x & 2.305\\
10200 & 512 & 256 & 256 & 39.124 & 133.1 & 6.985 & 64.2 & 0.48x & 3.740\\
10200 & 1024 & 64 & 64 & 36.046 & 144.5 & 14.237 & 63.0 & 0.44x & 1.942\\
10200 & 1024 & 128 & 128 & 34.395 & 151.4 & 11.835 & 75.7 & 0.50x & 1.734\\
10200 & 1024 & 256 & 256 & 33.747 & 154.3 & 11.066 & 81.0 & 0.52x & 2.155\\
10200 & 4096 & 64 & 64 & 33.471 & 155.6 & 46.471 & 77.1 & 0.50x & 5.605\\
10200 & 4096 & 128 & 128 & 33.413 & 155.9 & 42.642 & 84.1 & 0.54x & 5.139\\
10200 & 4096 & 256 & 256 & 33.398 & 155.9 & 41.080 & 87.3 & 0.56x & 4.471\\
50624 & 8 & 64 & 8 & 121.100 & 213.4 & 0.086 & 81.5 & 0.38x & 0.117\\
50624 & 8 & 128 & 8 & 119.657 & 216.0 & 0.068 & 102.7 & 0.48x & 0.101\\
50624 & 8 & 256 & 8 & 119.683 & 216.0 & 0.066 & 105.8 & 0.49x & 0.034\\
50624 & 64 & 64 & 64 & 119.940 & 215.5 & 2.229 & 25.1 & 0.12x & 0.479\\
50624 & 64 & 128 & 64 & 119.985 & 215.4 & 2.126 & 26.3 & 0.12x & 0.475\\
50624 & 64 & 256 & 64 & 119.686 & 216.0 & 2.172 & 25.8 & 0.12x & 0.492\\
50624 & 512 & 64 & 64 & 151.092 & 171.1 & 9.999 & 44.8 & 0.26x & 1.690\\
50624 & 512 & 128 & 128 & 154.407 & 167.4 & 8.189 & 54.7 & 0.33x & 2.159\\
50624 & 512 & 256 & 256 & 155.417 & 166.3 & 7.165 & 62.5 & 0.38x & 2.081\\
50624 & 1024 & 64 & 64 & 150.644 & 171.6 & 10.847 & 82.6 & 0.48x & 2.664\\
50624 & 1024 & 128 & 128 & 140.629 & 183.8 & 11.602 & 77.3 & 0.42x & 3.053\\
50624 & 1024 & 256 & 256 & 123.319 & 209.6 & 9.394 & 95.4 & 0.46x & 3.803\\
50624 & 4096 & 64 & 64 & 139.898 & 184.8 & 38.646 & 92.8 & 0.50x & 7.171\\
50624 & 4096 & 128 & 128 & 151.190 & 171.0 & 34.846 & 102.9 & 0.60x & 7.514\\
50624 & 4096 & 256 & 256 & 129.098 & 200.2 & 35.584 & 100.8 & 0.50x & 8.489\\
\end{longtable}
}

{\tiny\setlength{\tabcolsep}{2.3pt}
\begin{longtable}{rrrrrrrrrr}
\caption{100k 网格（94378 cells）：SI/RSI 并行效率；缺少 S224/4096/请求 B=256}\label{tab:factor-efficiency-100k}\\
\toprule
方向数 & 样本数 & 请求 B & 实际 B & SI sweep/s & $E_{\rm SI}$ & RSI sweep/s & $E_{\rm RSI}$ & RSI/SI & RSI plan/s\\
\midrule
\endfirsthead
\multicolumn{10}{l}{\small 续表：100k 网格（94378 cells）：SI/RSI 并行效率；缺少 S224/4096/请求 B=256}\\
\toprule
方向数 & 样本数 & 请求 B & 实际 B & SI sweep/s & $E_{\rm SI}$ & RSI sweep/s & $E_{\rm RSI}$ & RSI/SI & RSI plan/s\\
\midrule
\endhead
\bottomrule
\endfoot
120 & 8 & 64 & 8 & 0.086 & 1845.5 & 0.148 & 122.2 & 0.07x & 0.297\\
120 & 8 & 128 & 8 & 0.092 & 1728.9 & 0.144 & 126.2 & 0.07x & 0.063\\
120 & 8 & 256 & 8 & 0.085 & 1857.0 & 0.131 & 138.1 & 0.07x & 0.064\\
120 & 64 & 64 & 64 & 0.086 & 1838.1 & 0.173 & 839.7 & 0.46x & 0.035\\
120 & 64 & 128 & 64 & 0.085 & 1858.2 & 0.159 & 910.9 & 0.49x & 0.035\\
120 & 64 & 256 & 64 & 0.085 & 1870.9 & 0.162 & 894.9 & 0.48x & 0.035\\
120 & 512 & 64 & 64 & 0.086 & 1845.2 & 1.235 & 939.1 & 0.51x & 0.138\\
120 & 512 & 128 & 128 & 0.084 & 1882.5 & 1.081 & 1073.2 & 0.57x & 0.078\\
120 & 512 & 256 & 256 & 0.085 & 1875.8 & 0.826 & 1403.5 & 0.75x & 0.049\\
120 & 1024 & 64 & 64 & 0.084 & 1887.2 & 2.780 & 834.4 & 0.44x & 0.256\\
120 & 1024 & 128 & 128 & 0.085 & 1872.9 & 2.118 & 1094.9 & 0.58x & 0.138\\
120 & 1024 & 256 & 256 & 0.084 & 1896.1 & 1.682 & 1378.7 & 0.73x & 0.078\\
120 & 4096 & 64 & 64 & 0.085 & 1873.0 & 11.600 & 799.8 & 0.43x & 0.980\\
120 & 4096 & 128 & 128 & 0.085 & 1873.1 & 8.351 & 1111.0 & 0.59x & 0.496\\
120 & 4096 & 256 & 256 & 0.092 & 1714.2 & 6.849 & 1354.5 & 0.79x & 0.259\\
440 & 8 & 64 & 8 & 0.316 & 1837.6 & 0.147 & 122.9 & 0.07x & 0.390\\
440 & 8 & 128 & 8 & 0.325 & 1790.6 & 0.138 & 131.7 & 0.07x & 0.098\\
440 & 8 & 256 & 8 & 0.318 & 1828.5 & 0.252 & 72.0 & 0.04x & 0.097\\
440 & 64 & 64 & 64 & 0.339 & 1714.2 & 0.160 & 904.1 & 0.53x & 1.086\\
440 & 64 & 128 & 64 & 0.341 & 1706.6 & 0.168 & 862.0 & 0.51x & 0.234\\
440 & 64 & 256 & 64 & 0.358 & 1625.5 & 0.151 & 959.7 & 0.59x & 0.239\\
440 & 512 & 64 & 64 & 0.318 & 1829.4 & 1.247 & 930.2 & 0.51x & 0.666\\
440 & 512 & 128 & 128 & 0.317 & 1833.4 & 1.021 & 1136.3 & 0.62x & 0.396\\
440 & 512 & 256 & 256 & 0.323 & 1798.6 & 0.850 & 1364.1 & 0.76x & 0.289\\
440 & 1024 & 64 & 64 & 0.321 & 1811.3 & 3.134 & 740.0 & 0.41x & 1.028\\
440 & 1024 & 128 & 128 & 0.386 & 1507.3 & 2.282 & 1016.5 & 0.67x & 0.715\\
440 & 1024 & 256 & 256 & 0.389 & 1495.3 & 2.051 & 1131.0 & 0.76x & 0.433\\
440 & 4096 & 64 & 64 & 0.377 & 1543.5 & 12.316 & 753.3 & 0.49x & 3.945\\
440 & 4096 & 128 & 128 & 0.334 & 1739.8 & 7.908 & 1173.2 & 0.67x & 2.107\\
440 & 4096 & 256 & 256 & 0.323 & 1797.3 & 7.059 & 1314.4 & 0.73x & 1.025\\
960 & 8 & 64 & 8 & 0.805 & 1575.2 & 0.157 & 115.1 & 0.07x & 0.550\\
960 & 8 & 128 & 8 & 0.809 & 1568.7 & 0.158 & 114.5 & 0.07x & 0.125\\
960 & 8 & 256 & 8 & 0.822 & 1542.2 & 0.171 & 105.9 & 0.07x & 0.126\\
960 & 64 & 64 & 64 & 0.819 & 1548.6 & 0.194 & 748.1 & 0.48x & 2.630\\
960 & 64 & 128 & 64 & 0.823 & 1541.4 & 0.185 & 781.8 & 0.51x & 0.580\\
960 & 64 & 256 & 64 & 0.814 & 1558.6 & 0.161 & 903.0 & 0.58x & 0.620\\
960 & 512 & 64 & 64 & 0.814 & 1557.6 & 1.517 & 764.6 & 0.49x & 1.854\\
960 & 512 & 128 & 128 & 0.823 & 1540.9 & 1.087 & 1066.7 & 0.69x & 1.107\\
960 & 512 & 256 & 256 & 0.816 & 1555.1 & 1.033 & 1123.1 & 0.72x & 0.754\\
960 & 1024 & 64 & 64 & 0.810 & 1565.2 & 2.840 & 816.6 & 0.52x & 2.322\\
960 & 1024 & 128 & 128 & 0.804 & 1577.9 & 2.131 & 1088.3 & 0.69x & 1.804\\
960 & 1024 & 256 & 256 & 0.809 & 1568.0 & 2.031 & 1141.8 & 0.73x & 0.985\\
960 & 4096 & 64 & 64 & 0.819 & 1548.6 & 11.486 & 807.8 & 0.52x & 7.736\\
960 & 4096 & 128 & 128 & 0.734 & 1728.2 & 7.703 & 1204.4 & 0.70x & 5.540\\
960 & 4096 & 256 & 256 & 0.709 & 1788.4 & 7.067 & 1312.9 & 0.73x & 2.247\\
5040 & 8 & 64 & 8 & 3.656 & 1821.5 & 0.120 & 150.6 & 0.08x & 0.532\\
5040 & 8 & 128 & 8 & 3.675 & 1812.0 & 0.130 & 139.9 & 0.08x & 0.077\\
5040 & 8 & 256 & 8 & 3.613 & 1843.3 & 0.203 & 89.3 & 0.05x & 0.077\\
5040 & 64 & 64 & 64 & 3.597 & 1851.6 & 0.189 & 768.8 & 0.42x & 3.740\\
5040 & 64 & 128 & 64 & 3.608 & 1845.6 & 0.222 & 653.6 & 0.35x & 1.208\\
5040 & 64 & 256 & 64 & 3.658 & 1820.6 & 0.188 & 769.5 & 0.42x & 1.336\\
5040 & 512 & 64 & 64 & 3.675 & 1812.1 & 1.354 & 856.7 & 0.47x & 10.947\\
5040 & 512 & 128 & 128 & 3.610 & 1844.6 & 0.934 & 1242.3 & 0.67x & 3.673\\
5040 & 512 & 256 & 256 & 3.960 & 1681.6 & 0.919 & 1262.4 & 0.75x & 3.805\\
5040 & 1024 & 64 & 64 & 3.679 & 1810.2 & 2.893 & 801.8 & 0.44x & 7.350\\
5040 & 1024 & 128 & 128 & 3.691 & 1804.1 & 2.112 & 1098.4 & 0.61x & 5.125\\
5040 & 1024 & 256 & 256 & 4.253 & 1566.0 & 1.925 & 1204.7 & 0.77x & 5.215\\
5040 & 4096 & 64 & 64 & 3.633 & 1833.1 & 11.088 & 836.7 & 0.46x & 17.224\\
5040 & 4096 & 128 & 128 & 3.611 & 1844.0 & 8.933 & 1038.6 & 0.56x & 19.105\\
5040 & 4096 & 256 & 256 & 4.266 & 1561.1 & 8.244 & 1125.4 & 0.72x & 14.579\\
10200 & 8 & 64 & 8 & 8.400 & 1604.4 & 0.143 & 126.4 & 0.08x & 0.619\\
10200 & 8 & 128 & 8 & 8.636 & 1560.6 & 0.150 & 121.1 & 0.08x & 0.093\\
10200 & 8 & 256 & 8 & 8.631 & 1561.4 & 0.158 & 114.8 & 0.07x & 0.101\\
10200 & 64 & 64 & 64 & 8.037 & 1676.8 & 0.142 & 1023.8 & 0.61x & 1.007\\
10200 & 64 & 128 & 64 & 7.516 & 1793.2 & 0.183 & 793.6 & 0.44x & 1.074\\
10200 & 64 & 256 & 64 & 7.444 & 1810.4 & 0.197 & 737.3 & 0.41x & 1.028\\
10200 & 512 & 64 & 64 & 7.337 & 1836.8 & 1.290 & 898.7 & 0.49x & 3.765\\
10200 & 512 & 128 & 128 & 7.496 & 1798.0 & 0.954 & 1216.0 & 0.68x & 3.798\\
10200 & 512 & 256 & 256 & 7.302 & 1845.8 & 0.936 & 1238.6 & 0.67x & 6.227\\
10200 & 1024 & 64 & 64 & 8.740 & 1542.1 & 2.818 & 823.1 & 0.53x & 8.102\\
10200 & 1024 & 128 & 128 & 8.330 & 1617.9 & 1.923 & 1206.1 & 0.75x & 8.958\\
10200 & 1024 & 256 & 256 & 7.309 & 1844.0 & 1.753 & 1322.9 & 0.72x & 6.913\\
10200 & 4096 & 64 & 64 & 7.333 & 1838.0 & 11.155 & 831.7 & 0.45x & 21.271\\
10200 & 4096 & 128 & 128 & 7.418 & 1816.8 & 8.066 & 1150.2 & 0.63x & 23.987\\
10200 & 4096 & 256 & 256 & 7.391 & 1823.4 & 6.952 & 1334.5 & 0.73x & 21.259\\
50624 & 8 & 64 & 8 & 39.909 & 1676.0 & 0.142 & 128.0 & 0.08x & 0.656\\
50624 & 8 & 128 & 8 & 40.256 & 1661.6 & 0.143 & 127.0 & 0.08x & 0.107\\
50624 & 8 & 256 & 8 & 37.445 & 1786.3 & 0.142 & 127.3 & 0.07x & 0.096\\
50624 & 64 & 64 & 64 & 51.375 & 1302.0 & 0.174 & 834.3 & 0.64x & 6.143\\
50624 & 64 & 128 & 64 & 42.930 & 1558.1 & 0.275 & 527.8 & 0.34x & 1.108\\
50624 & 64 & 256 & 64 & 38.965 & 1716.7 & 0.154 & 941.2 & 0.55x & 0.866\\
50624 & 512 & 64 & 64 & 37.517 & 1782.9 & 1.451 & 799.4 & 0.45x & 4.026\\
50624 & 512 & 128 & 128 & 39.635 & 1687.6 & 1.085 & 1069.0 & 0.63x & 4.799\\
50624 & 512 & 256 & 256 & 41.169 & 1624.8 & 1.019 & 1137.9 & 0.70x & 6.587\\
50624 & 1024 & 64 & 64 & 38.720 & 1727.5 & 2.580 & 898.9 & 0.52x & 7.079\\
50624 & 1024 & 128 & 128 & 42.669 & 1567.6 & 2.140 & 1083.7 & 0.69x & 9.039\\
50624 & 1024 & 256 & 256 & 42.633 & 1569.0 & 2.050 & 1131.4 & 0.72x & 10.793\\
50624 & 4096 & 64 & 64 & 39.466 & 1694.9 & 11.997 & 773.3 & 0.46x & 27.635\\
50624 & 4096 & 128 & 128 & 42.007 & 1592.3 & 8.642 & 1073.6 & 0.67x & 33.531\\
\end{longtable}
}


\paragraph{影响因素与低误差最优组合}

端到端加速的首要影响因素是角度方向数：SI 必须计算全部 $N_{\Omega}$ 个方向，而 RSI 每步只推进一条样本链。固定 30k 网格、4,096 样本和请求 $B=128$ 时，$S_{70}$（5,040 directions）、$S_{100}$（10,200 directions）和 $S_{224}$（50,624 directions）的端到端加速比分别为 $0.345\times$、$0.723\times$ 和 $3.740\times$；相应 error 为 $1.945\%$、$1.602\%$ 和 $1.054\%$。样本数是精度--时间权衡的主导因素：同一 30k/$S_{224}$/$B=128$ 下，8、64、512、1,024、4,096 样本的 error 从 $35.121\%$、$19.619\%$、$4.767\%$、$2.395\%$ 降至 $1.054\%$，但端到端加速也从 $759.439\times$、$48.210\times$、$15.456\times$、$10.040\times$ 降至 $3.740\times$。batch 则主要影响 RSI 的 kernel 吞吐和 plan 调度；在该低误差点，$B=64/128/256$ 的端到端加速分别为 $3.179\times$、$3.740\times$、$3.111\times$，故 $B=128$ 为本机的甜点点。空间网格会进一步改变 plan/cache 成本；例如同一 $S_{224}$/4,096/$B=128$ 在 100k 上的单 sweep 吞吐比为 $0.674\times$，高于 30k 的 $0.602\times$，但 100k 的 RSI plan 时间为 33.531 s，端到端加速降为 $3.112\times$。

在 error 不超过 $2\%$ 的 23 组中，30k/$S_{224}$/4,096/$B=128$ 的端到端加速最高，且 error 在全部 269 组中最低。该 case 的 SI 收敛轮数为 14，RSI-tail 使用 $14+10=24$ 步，因此 SI 需要 $50,624\times14=708,736$ 次 directional cell-sweep/cell，RSI 需要 $4,096\times24=98,304$ 次 sample-sweep/cell，工作量减少 $7.21\times$。其 sweep kernel 统计如下；block 和 logical thread 均为整个 SI fine 或 RSI 过程的累计 launch 几何，不表示硬件实际活跃 lane 数。

{\scriptsize\setlength{\tabcolsep}{3.3pt}
\begin{tabular}{lrrrrrr}
\toprule
路径 & launches & blocks & logical threads & threads/block & $E$ (Mcell-sweeps/s) & $E_{\rm RSI}/E_{\rm SI}$\\
\midrule
SI & 7,602 & 710,500 & 181,637,120 & 256 & 170.96 & --\\
RSI & 209,666 & 80,314,144 & 20,560,154,624 & 256 & 102.88 & 0.602$\times$\\
\bottomrule
\end{tabular}}

因此该 case 的 sweep 阶段实际加速为 $151.190/34.846=4.338\times$；RSI 的单位 sweep 吞吐只有 SI 的 $0.602\times$，端到端的 $3.740\times$ 优势来自所需方向/样本 sweep 数显著减少，而不是 RSI 单个 kernel 的并行吞吐更高。

\paragraph{Nsight Systems/Compute 采样}

对同一 case 的完整 Nsight Systems trace 共记录 227,160 次 \code{cudaLaunchKernel}。GPU kernel 时间的 $62.9\%$ 位于 \code{sweepSamplesKernel}，$36.5\%$ 位于 \code{sweepLevelFusedKernel}。\\
主机端 kernel-launch API 仅占 $0.3\%$ 时间；其余等待主要对应 GPU 同步，而非 host launch 调用。该 case 的主问题因此是 GPU sweep 执行。trace 下的绝对时间受 profiling 记录开销影响，以下仅使用其 kernel 占比和实例数作归因。

{\scriptsize\setlength{\tabcolsep}{3.3pt}
\begin{tabular}{lrrr}
\toprule
Nsight Systems kernel & GPU time share & instances & total GPU time/s\\
\midrule
\code{sweepSamplesKernel} & 62.9\% & 560 & 131.323\\
\code{sweepLevelFusedKernel} & 36.5\% & 1,281 & 76.209\\
\bottomrule
\end{tabular}}

Nsight Compute 对两个主导 kernel 各采一个代表性 launch；这些硬件计数用于衡量实际活跃 warp、缓存层级和 launch 尺度，不能由累计 block 数替代。\code{sweepSamplesKernel} 与 \code{sweepLevelFusedKernel} 都由 SI fine 和 RSI 路径复用，现有 kernel-name 采样没有 NVTX 方法标签，故不能把单次采样严格归属给 SI 或 RSI；SI 专属判断以下面的完整运行计时、plan cache 和 launch 几何为准。\code{sweepSamplesKernel} 的采样 grid 只有 1 个 block，表明某些 cyclic/fallback launch 存在严重的单 block 低占用；\code{sweepLevelFusedKernel} 的 grid 为 128 个 block，仍只有 $0.37$ waves/SM，且其 L1 命中率低、L2 命中率高，符合 L2/访存和 wavefront 负载受限的特征。

{\scriptsize\setlength{\tabcolsep}{2.8pt}
\begin{tabular}{lrrrrrrrr}
\toprule
NCU sampled kernel & grid & block & achieved/theoretical occ. & waves/SM & L1 hit & L2 hit & DRAM throughput\\
\midrule
\code{sweepSamplesKernel} & 1 & 128 & 2.08\% / 83.33\% & 0.00 & 15.33\% & 10.34\% & 1.32\%\\
\code{sweepLevelFusedKernel} & 128 & 256 & 30.35\% / 83.33\% & 0.37 & 9.48\% & 97.14\% & 7.16\%\\
\bottomrule
\end{tabular}}

两次 NCU 采样还分别报告约 $39\%$ 和 $75\%$ 的 global-memory sectors 为 excessive/uncoalesced。优化优先级因此为：先减少 cyclic \code{sweepSamplesKernel} 的低 grid 规模或将其与其他方向/样本合并，再改善 SI/RSI 共用 fused wavefront 的全局访问合并度和 level 内负载均衡；单纯增加 block 数不能解决这两类瓶颈。若需把 occupancy 和缓存命中率精确拆给 SI 与 RSI，应增加 NVTX range 或 SI-only profile。

\paragraph{数据、含义与结论汇总}

下表将本节中容易混淆的计时、launch 几何和硬件计数统一解释。除 NCU 的两个代表性 launch 外，其余数据均来自 30k/$S_{224}$/4,096/$B=128$ 的完整运行；NCU 单次采样只用于定位瓶颈，不能外推为每个 kernel 实例的精确平均值，也不能仅据 kernel 名拆分为 SI 或 RSI。

{\scriptsize\renewcommand{\arraystretch}{1.10}\setlength{\tabcolsep}{3pt}
\begin{tabularx}{\linewidth}{>{\raggedright\arraybackslash}p{0.19\linewidth}>{\raggedright\arraybackslash}p{0.22\linewidth}>{\raggedright\arraybackslash}X>{\raggedright\arraybackslash}p{0.22\linewidth}}
\toprule
指标 & 数据 & 含义 & 结论\\
\midrule
低误差最优点 & error $=1.054\%$；端到端加速 $=3.740\times$ & 在 269 组中 error 最低，且在 error $\leq2\%$ 的 23 组中端到端加速最高。 & 30k/$S_{224}$/4,096/$B=128$ 是当前精度约束下的推荐配置。\\
工作量与单位吞吐 & SI/RSI sweep 工作量 $=708,736/98,304=7.21\times$；$E_{\rm RSI}/E_{\rm SI}=0.602\times$ & RSI 每个等价 cell-sweep 的吞吐低于 SI；但待执行的样本 sweep 数少得多。 & $3.740\times$ 加速主要来自工作量减少，不表示 RSI 单位计算更快。\\
累计 launch 几何 & SI：7,602 launches、710,500 blocks；RSI：209,666 launches、80,314,144 blocks；均为 256 threads/block & 这是完整计算中各 kernel grid 的累加，包含方向、level、样本 batch 和 cyclic 回退的不同分解。 & RSI block 总数大于 SI 并不构成比较不公平；它不能衡量实际 SM 活跃线程或硬件利用率。\\
SI sweep 时间与吞吐 & $N_\Omega=50,624$、$N=14$、$t_{\rm SI,sweep}=151.190$ s，占 $t_{\rm SI,total}=168.836$ s 的 $89.55\%$；$E_{\rm SI}=170.96$ Mcell-sweeps/s & SI 已完成 $708,736$ directional cell-sweeps/cell；total 时间主要由 GPU sweep 决定，norm/reduce/copy 不构成主导。 & SI 优化应首先降低波前 sweep 时间，例如提高同一 level 内并行度和访问合并度；不应把重点放在尾部归约。\\
SI plan/cache 与 launch 形状 & 396 个方向 plan chunks；device cache hits/misses $=5,148/396$；首次读取/上传 $=6.925/6.923$ GiB，速率 $=0.760/2.561$ GiB/s；7,602 launches、710,500 blocks，即平均 93.46 blocks/launch & 方向 plan 首次物化与 H2D 传输是 SI 的固定成本；累计 block 多数来自方向块与波前层次，平均 grid 仍不足以推断每次 launch 的占用率。 & 保持 device plan cache，并优先重叠 plan 物化/传输与 sweep；SI 的真实 occupancy 需要 SI-only 或 NVTX 标记 NCU，不能由平均 blocks/launch 替代。\\
主导 GPU 时间 & \code{sweepSamplesKernel}：$62.9\%$、560 instances；\code{sweepLevelFusedKernel}：$36.5\%$、1,281 instances；host launch API：$0.3\%$ & 时间主要花在 GPU sweep，本机 host 发射开销不是主导项。 & 优化应优先针对两个 sweep kernel 的并行度和访存，而非减少 host API 调用。\\
cyclic/fallback 代表性利用率 & \code{sweepSamplesKernel}: grid $=1$、block $=128$、achieved occupancy $=2.08\%$（理论 $83.33\%$） & 该代表性 launch 只有一个 block，远不足以覆盖 GPU 的 SM；因 kernel 为共用实现，不能仅据此断定来源为 SI 或 RSI。 & 优先合并或批处理 cyclic/fallback 工作，减少单 block launch；仅增大每 block 线程数无效。\\
共用 fused wavefront 代表性利用率与缓存 & \code{sweepLevelFusedKernel}: grid $=128$、block $=256$、occupancy $=30.35\%$、waves/SM $=0.37$、L1/L2 hit $=9.48\%/97.14\%$、DRAM $=7.16\%$ peak、uncoalesced $\approx75\%$ & wavefront 已有一定并行度，但 level 宽度不足、L1 复用低且全局访问合并度差；高 L2 hit 不等于高 DRAM 带宽利用。 & 这同时是 SI 与 RSI 的优化候选：改善 level 内负载均衡和数据布局/访问合并度，而不是依据累计 block 数继续加大 grid。\\
\bottomrule
\end{tabularx}}

\sect{9. 关键可调参数}

\begin{tabularx}{\linewidth}{>{\ttfamily}l X}
\toprule
环境变量 & 作用\\
\midrule
RSI\_CUDA\_STREAMING\_PLAN\_THRESHOLD\_MB & 控制何时从全量 plan 切换到流式 sweep plan。\\
RSI\_CUDA\_SI\_PLAN\_CHUNK & SI 每个 plan 方向块的方向数，默认 128。\\
RSI\_CUDA\_SI\_DEVICE\_PLAN\_CACHE\_MB & SI 设备端 plan 常驻显存预算。\\
RSI\_CUDA\_PLAN\_HOST\_CACHE\_MB & 主机端 LRU plan 缓存预算。\\
RSI\_SWEEP\_PLAN\_CACHE\_DIR & 磁盘 plan 缓存目录。\\
RSI\_CUDA\_SI\_WAVEFRONT & 开关 SI 无环波前并行扫掠。\\
RSI\_CUDA\_SI\_TILED\_WAVEFRONT & 开关 SI 依赖层分块并行。\\
RSI\_CUDA\_RSI\_WAVEFRONT & 开关 RSI 无环波前并行扫掠。\\
RSI\_CUDA\_RSI\_TILED\_WAVEFRONT & 开关 RSI 依赖层分块并行。\\
RSI\_CUDA\_RSI\_FUSED\_WAVEFRONT & 开关 RSI fused wavefront；大网格默认开启。\\
RSI\_CUDA\_RSI\_MAX\_SAMPLES\_PER\_BATCH & 控制 streaming RSI 每个 batch 的样本数，默认 128；host memory 紧张时可调小。\\
RSI\_CUDA\_RSI\_PLAN\_PREFETCH & 是否在 GPU sweep 当前 RSI batch 时预取下一 batch 的 host plan，默认开启。\\
RSI\_CUDA\_RSI\_ADAPTIVE\_WAVEFRONT & 是否启用 fused/tiled 自适应实验路径，默认关闭。\\
RSI\_CUDA\_RSI\_ADAPTIVE\_TILE\_THRESHOLD & 自适应路径切回 tiled wavefront 的 level tile 阈值。\\
RSI\_CUDA\_RSI\_GRAPH\_WAVEFRONT & 是否对 streaming RSI level-by-level 路径做 CUDA Graph 捕获实验，默认关闭。\\
RSI\_CUDA\_RSI\_FACE\_LANE\_WAVEFRONT & 是否启用 face-lane wavefront 实验 kernel，默认关闭。\\
RSI\_CUDA\_FIXED\_PLAN\_CHUNK & RSI 固定方向缓存的方向粒度，默认 1。\\
RSI\_CUDA\_FIXED\_PLAN\_PARALLEL\_LOAD & 是否并行加载固定方向 plan 缓存。\\
RSI\_CUDA\_RSI\_PLAN\_PREWARM & 是否在 RSI 前预热固定方向缓存。\\
RSI\_CUDA\_FIXED\_PLAN\_PREWARM\_WORKERS & 预热固定方向缓存的 CPU worker 数。\\
RSI\_CUDA\_FIXED\_PLAN\_MISS\_DEDUP & 是否对固定方向缓存 miss 去重。\\
\bottomrule
\end{tabularx}

\sect{10. 当前优化逻辑总结}

当前实现的核心原则是：

\begin{enumerate}
\item SI 用角度并行和波前并行，把 ``方向间独立'' 和 ``同一依赖层内单元独立'' 都利用起来。
\item RSI 用样本并行和样本批内紧凑 plan，只为当前样本批真正用到的方向上传 plan。
\item 对 230k + S128 这类超大 plan，避免全量常驻，改为流式 sweep plan。
\item SI 的方向访问重复且有序，因此优先做设备端 plan 缓存。
\item RSI 的方向访问随机且样本批变化，因此优先做固定方向缓存、预热和并行加载。
\item RSI 的无环 sweep 在大网格上默认使用 fused wavefront，减少逐 level kernel launch，并保持和原 level-by-level wavefront 相同的依赖顺序。
\item RSI 与 RSI-tail 合并在同一批样本链中计算，避免重复随机链和重复 sweep。
\item mesh 数据按单元-面引用展平为 SoA，减少扫掠 kernel 中的间接访存。
\item 收敛范数、方向累加、样本 reduce 都在 GPU 上完成，只把最终场和极小控制量拷回 CPU。
\end{enumerate}

整体上，SI 的优化重点是 ``少存完整角通量、少重复上传 plan、在无环依赖层内打开单元并行''；RSI 的优化重点是 ``样本批处理、只上传实际方向、复用方向 plan、把 RSI 和 tail 合并、减少大网格逐 level launch''。

\end{document}
